% ══════════════════════════════════════════════════════════════════════════════
%  DATect Paper — MDPI Toxins Format
%  Converted from elsarticle format with expanded content
% ══════════════════════════════════════════════════════════════════════════════
%
% METRICS UPDATED: Post-infrastructure overhaul (monotonic constraints, Oregon feature reduction).
% Standard eval (seed=123, 40%, N=2,181): Ensemble R²=0.215, MAE=6.42. CIs pending scripts/eval/eval_paper_metrics.py re-run.
% Temporal holdout (2019-2023, N=1,036): Pooled R²=0.315, WA mean=0.549, OR mean=0.113.
% Spike metrics (hybrid alert): transition recall=0.734, event recall=0.859 (up from 0.652/0.815).
% Ablation table uses development set (seed=42, N=1,177) as intended; values unchanged.
% Dev set reference (seed=42): R²=0.414, MAE=6.380, RMSE=16.434.
% Bootstrap CIs in Tables 1/2 are final values from scripts/eval/eval_paper_metrics.py (B=10,000, seed=123, N=2,181).
%

\documentclass[toxins,article,submit,pdftex,moreauthors]{Definitions/mdpi}

% ── Additional Packages ─────────────────────────────────────────────────────
\usepackage{amsmath,amssymb}
\usepackage{siunitx}
\usepackage[ruled,vlined]{algorithm2e}
\usepackage{enumitem}

% ── SI units ─────────────────────────────────────────────────────────────────
\DeclareSIUnit{\microgrampeg}{\micro g\per g}
\DeclareSIUnit{\cellsperL}{cells\per L}

% ── MDPI internal commands ───────────────────────────────────────────────────
\firstpage{1}
\makeatletter
\setcounter{page}{\@firstpage}
\makeatother
\pubvolume{1}
\issuenum{1}
\articlenumber{0}
\pubyear{2026}
\copyrightyear{2026}
\datereceived{ }
\daterevised{ }
\dateaccepted{ }
\datepublished{ }
\hreflink{https://doi.org/}

% ══════════════════════════════════════════════════════════════════════════════
%  METADATA
% ══════════════════════════════════════════════════════════════════════════════

\Title{DATect: Per-Site Machine Learning Forecasting of Domoic Acid Along the Pacific Northwest Coast}

\TitleCitation{DATect: Per-Site ML Forecasting of Domoic Acid in the Pacific Northwest}

\newcommand{\orcidauthorA}{0000-0000-0000-000X}

\Author{Anson Chen $^{1,*}$\orcidA{} and John Mickett $^{2}$}

\AuthorNames{Anson Chen and John Mickett}

\AuthorCitation{Chen, A.; Mickett, J.}

\address{%
$^{1}$ \quad Paul G. Allen School of Computer Science \& Engineering, University of Washington, Seattle, WA 98195, USA; ac283@uw.edu\\
$^{2}$ \quad Applied Physics Laboratory, University of Washington, Seattle, WA 98105, USA; mickett@uw.edu}

\corres{Correspondence: ac283@uw.edu}

% ══════════════════════════════════════════════════════════════════════════════
%  ABSTRACT
% ══════════════════════════════════════════════════════════════════════════════
\abstract{
Domoic acid (DA) contamination of shellfish is a recurring public health crisis along the U.S.\ Pacific Northwest coast, yet no quantitative forecasting system exists for Washington and Oregon. We present DATect, a one-week-ahead forecasting system for 10 monitoring sites spanning \SI{600}{\kilo\meter} of coastline, trained on 21 years of satellite, climate, hydrological, and biological data (2003--2023). DATect uses per-site XGBoost or Random Forest models and is evaluated with leak-free retrospective validation on a held-out test set ($n = 2{,}181$) kept separate from configuration development. All five Washington sites achieve high skill ($R^2 = 0.480$--$0.789$), whereas Oregon sites show lower skill consistent with sparser and less regular monitoring. At the operational \SI{20}{\microgrampeg} threshold, thresholding the regression ensemble detects only 12.4\% of held-out below-to-above-threshold transitions, while na\"ive persistence detects 23.6\%. A dedicated binary classifier trained with a safe-baseline filter detects 73.4\% of these transitions, a 3.1$\times$ improvement over na\"ive persistence ($0.734 / 0.236$). These results show that actionable early warning requires a transition-focused detector in addition to magnitude forecasting. DATect is deployed as a publicly accessible web dashboard.
}

\keyword{domoic acid; harmful algal blooms; machine learning; ensemble forecasting; \textit{Pseudo-nitzschia}; Pacific Northwest; XGBoost; data leakage prevention}

% ══════════════════════════════════════════════════════════════════════════════
%  DOCUMENT
% ══════════════════════════════════════════════════════════════════════════════
\begin{document}

% ══════════════════════════════════════════════════════════════════════════════
%  1. INTRODUCTION
% ══════════════════════════════════════════════════════════════════════════════
\section{Introduction}
\label{sec:introduction}

% ── 1.1 DA Biology, Drivers, and Impacts ─────────────────────────────────────
\subsection{Domoic Acid and \textit{Pseudo-nitzschia} Blooms}
\label{sec:intro-biology}

Domoic acid (DA) is a potent neurotoxin produced by certain species of the marine diatom genus \textit{Pseudo-nitzschia} and is the causative agent of amnesic shellfish poisoning \cite{Bates2018, Trainer2012, Wright1989}. Since the first documented outbreak in 1987, DA contamination has become a persistent threat along the U.S.\ Pacific coast, with severe consequences for human health, coastal economies, and marine wildlife \cite{Mos2001, Lefebvre2010, Trainer2012}. The 2015 marine heatwave illustrates the scale of this risk: DA concentrations in razor clams exceeded the federal safety limit by over 700\%, causing tens of millions of dollars in fishery losses \cite{McCabe2016, Bond2015, Trainer2020}.

Forecasting DA is difficult because toxin production is only weakly coupled to bloom magnitude. DA production depends strongly on physiological stress---particularly nutrient limitation---so a large \textit{Pseudo-nitzschia} bloom may be weakly toxic while a smaller bloom generates dangerous contamination \cite{Trick2010, Marchetti2004, Lelong2012}. In the Pacific Northwest (PNW), this biology interacts with coastal upwelling, climate variability, and retention zones such as the Juan de Fuca eddy and Heceta Bank, creating a multiscale nonlinear system with strong spatial heterogeneity \cite{Trainer2002, Trainer2009, McKibben2017, Jacox2018}.

The forecasting problem is further complicated by the monitoring data themselves: DA measurements are sparse, irregularly sampled, and site-dependent, while the target distribution is dominated by zeros punctuated by rare but operationally critical spikes. These properties motivate both DATect's feature design and its emphasis on leak-free evaluation.

% ── 1.2 Existing Monitoring and Forecasting ──────────────────────────────────
\subsection{Existing Monitoring and Forecasting in the Pacific Northwest}
\label{sec:intro-monitoring}

Current DA risk management in the PNW is reactive: beaches are closed only after measured toxin levels exceed the FDA regulatory limit of \SI{20}{\microgrampeg} \cite{McCabe2023}, giving managers no advance warning. The primary forecasting tool is the NOAA Pacific Northwest HAB Bulletin \cite{McCabe2023}, issued biweekly to monthly during peak bloom season. The Bulletin synthesizes environmental and biological data into qualitative risk categories (low/medium/high) that provide strategic planning value \cite{Jin2024} but do not produce numerical DA predictions and operate at seasonal rather than weekly resolution. Complementary regional resources---Oregon's phytoplankton monitoring program \cite{McKibben2015}, the SoundToxins network \cite{Trainer2023}, and the LiveOcean circulation model \cite{MacCready2021}---extend observational coverage but do not forecast DA. The California C-HARM system \cite{Anderson2016} provides bloom-probability nowcasts for California but does not extend to Washington or Oregon (Table~\ref{tab:system-comparison}).

A critical gap thus exists: \textbf{no quantitative DA concentration forecasting system has been developed for the Pacific Northwest}. All existing monitoring is reactive, and the sole forecasting tool provides qualitative risk categories, not numerical toxin predictions.

% ── 1.3 Problem Statement ────────────────────────────────────────────────────
\subsection{Problem Statement and Contributions}
\label{sec:intro-problem}

This paper addresses the absence of quantitative DA forecasting in the PNW by developing DATect, the first system to produce site-specific, numerical DA concentration predictions for this region. We formalize the prediction task as follows: given environmental features available at time $t - 7$ days (the ``anchor date''), predict the DA concentration (\si{\microgrampeg}) at time $t$ for a specific monitoring site $s$ among 10 Pacific Coast locations.

Two evaluation targets matter operationally. The first is magnitude forecasting: how accurately can a system predict DA concentration one week ahead? The second is transition detection: can the system anticipate below-to-above-threshold crossings at \SI{20}{\microgrampeg}, the events that trigger harvest closures? DATect is evaluated on both targets because a model can achieve useful average regression error while still missing the closure-relevant transitions that matter most to managers.

This task presents four compounding challenges:
\begin{enumerate}[nosep]
    \item \textbf{Sparse, irregular target variable}: DA measurements are collected approximately weekly but with frequent gaps, missingness, and inconsistent timing across sites.
    \item \textbf{High spatial variability}: The 10 monitoring sites span 5\textdegree\ of latitude ($\sim$\SI{600}{\kilo\meter}), with dramatically different oceanographic regimes, bloom histories, and DA dynamics.
    \item \textbf{Extreme value distribution}: DA measurements are dominated by zeros and low values, with rare but consequential spikes exceeding \SI{100}{\microgrampeg}.
    \item \textbf{Strict temporal integrity requirements}: Operational forecasting demands that no future information leak into predictions, requiring rigorous validation protocols.
\end{enumerate}

This paper makes four main contributions:
\begin{enumerate}[nosep]
    \item \textbf{Quantitative PNW DA forecasting under leak-free evaluation}: DATect provides the first site-specific, one-week-ahead DA concentration forecasts for Washington and Oregon, using 21 years of data across 10 monitoring sites and a held-out retrospective test set separated from configuration development.
    \item \textbf{Transition-focused early warning}: A dedicated binary classifier trained with a safe-baseline filter provides advance warning for below-to-above-threshold crossings that trigger harvest closures, complementing the regression forecast rather than replacing it (Section~\ref{sec:results-spike}).
    \item \textbf{Per-site modeling on sparse, irregular records}: Per-site XGBoost or Random Forest configurations, interpolated training, and raw-only persistence features provide a practical approach to sparse monitoring data while keeping evaluation tied to real observations.
    \item \textbf{Evidence that data quality is the binding constraint}: The Washington--Oregon performance split shows that monitoring density and regularity, not model complexity, are the primary limits on forecast skill in the current regional record.
\end{enumerate}


% ══════════════════════════════════════════════════════════════════════════════
%  2. RELATED WORK
% ══════════════════════════════════════════════════════════════════════════════
\section{Related Work}
\label{sec:related-work}

\subsection{Quantitative Forecasting of Domoic Acid and \textit{Pseudo-nitzschia}}
\label{sec:rw-da-forecasting}

Despite decades of DA monitoring and multiple qualitative forecasting systems, quantitative computational approaches specifically targeting DA or \textit{Pseudo-nitzschia} remain limited in number, geographic scope, and temporal extent \cite{Anderson2012, McGillicuddy2010}. No prior study has produced a quantitative DA concentration forecast for the PNW.

In the Galician R\'ias Baixas region of northwest Spain, Al\'aez et al.\ \cite{Alaez2021} applied Random Forest and AdaBoost classifiers to predict \textit{Pseudo-nitzschia} bloom events, achieving approximately 85\% classification accuracy. Their study covered approximately five monitoring stations over roughly a decade of data, targeting bloom occurrence (binary classification) rather than DA concentration. Cusack et al.\ \cite{Cusack2015} modeled \textit{Pseudo-nitzschia} events off southwest Ireland using environmental predictors, demonstrating the potential of empirical models for bloom forecasting in European waters. Additional studies have documented \textit{Pseudo-nitzschia} dynamics and DA contamination in new regions, including the Gulf of Maine \cite{Clark2021} and southern Florida \cite{Schreiber2023}, reflecting the global expansion of DA monitoring efforts. In the California Current, Lane et al.\ \cite{Lane2009} developed a logistic regression model for predicting toxigenic \textit{Pseudo-nitzschia} blooms in Monterey Bay using 8.3 years of monitoring data, achieving $\geq$75\% predictive accuracy with chlorophyll~$a$ and silicic acid as primary predictors.

The closest methodological parallel to DATect is the deep learning system developed by Grasso et al.\ \cite{Grasso2019} for forecasting paralytic shellfish toxin (PST) levels in coastal Maine. Using a convolutional neural network (CNN) trained on chemical data from routine monitoring sites, their system forecasts PST levels on a weekly basis with over 95\% prediction reliability within a one-week timeframe. However, this system targets a different toxin (PST from \textit{Alexandrium catenella}, not DA from \textit{Pseudo-nitzschia}), covers a smaller geographic area (approximately 10 sites along the Maine coast), and uses approximately five years of data.

The California Harmful Algae Risk Mapping (C-HARM) system \cite{Anderson2016} provides satellite-based nowcasts of \textit{Pseudo-nitzschia} bloom probability and DA risk along the California coast. C-HARM integrates satellite imagery with ocean model output to produce gridded risk maps but operates as a nowcast (current conditions) rather than a forecast and does not extend to Washington or Oregon. Moreno et al.\ \cite{Moreno2022} developed a mechanistic process-based model of \textit{Pseudo-nitzschia} and DA production for the California Current, representing a complementary approach to data-driven methods. Most recently, Brunson et al.\ \cite{Brunson2024} demonstrated molecular forecasting of DA using gene expression (dabA and sit1 coexpression) as a one-week-ahead predictor during a bloom in Monterey Bay, representing an emerging genomics-based approach that could eventually complement satellite- and ML-based systems like DATect.

\subsection{Machine Learning Methods for Harmful Algal Blooms}
\label{sec:rw-ml-habs}

Across HAB forecasting more broadly, tree ensembles, linear models, and deep sequence models have all been applied to heterogeneous environmental data \cite{Cruz2021, Shin2024, Hill2020, Yu2021, Manning2019}. For DATect, the most relevant lesson is not that one model family dominates universally, but that coastal forecasting methods must tolerate missingness, nonlinear interactions, and mixed satellite/climate/biological covariates. Gradient boosting and Random Forests therefore provide natural nonlinear competitors \cite{Friedman2001, Chen2016, Breiman2001}, while linear models remain useful baselines for testing whether the gains come from model flexibility or from feature engineering and site-specific design.

\subsection{Data Leakage in Environmental Machine Learning}
\label{sec:rw-leakage}

A growing body of work has highlighted the pervasive problem of data leakage in machine learning applications to environmental science. Kapoor and Narayanan \cite{Kapoor2023} demonstrated that data leakage affects a substantial fraction of published ML-based science papers, leading to inflated performance claims and irreproducible results. In temporal prediction tasks such as HAB forecasting, common sources of leakage include random train-test splits that violate temporal ordering, use of features computed from future data, and shared preprocessing across train and test sets \cite{Kaufman2012, Roberts2017, Bergmeir2012}. Proper out-of-sample evaluation for temporal data requires rolling-origin or expanding-window approaches rather than random holdouts \cite{Tashman2000}, and spatial autocorrelation in environmental data further complicates validation design \cite{Valavi2019}.

Many prior HAB forecasting studies employed random 80/20 train-test splits without temporal ordering \cite{Cruz2021, Shin2024}, which may have produced artificially inflated performance metrics. DATect addresses this explicitly through nine structural leakage-prevention mechanisms described in Section~\ref{sec:validation}.

\subsection{Comparison of Existing Systems}
\label{sec:rw-comparison}

\begin{table}[H]
\caption{Comparison of existing HAB forecasting systems with DATect. DATect provides the longest temporal record (21 years) and broadest geographic span ($\sim$\SI{600}{\kilo\meter}) of any published quantitative DA/\textit{Pseudo-nitzschia} ML forecasting study.}
\label{tab:system-comparison}
\small
\begin{tabular}{llllllll}
\toprule
\textbf{System} & \textbf{Region} & \textbf{Sites} & \textbf{Period} & \textbf{Method} & \textbf{Target} & \textbf{Horizon} & \textbf{Deployed} \\
\midrule
PNW HAB Bulletin & WA/OR & $\sim$20 & 2017-- & Qualitative & DA risk & 1--2 wk & Yes \\
\cite{McCabe2023} & & & & expert & (categorical) & & \\
\addlinespace
Maine PST & Maine & $\sim$10 & $\sim$5 yr & CNN & PST & 1 wk & No \\
\cite{Grasso2019} & & & & & (not DA) & & \\
\addlinespace
Galician RF & NW Spain & $\sim$5 & $\sim$10 yr & RF/AdaBoost & PN bloom & --- & No \\
\cite{Alaez2021} & & & & & (binary) & & \\
\addlinespace
C-HARM & California & Grid & Ongoing & Satellite+ & PN prob. & Nowcast & Yes \\
\cite{Anderson2016} & & & & model & & & \\
\addlinespace
Monterey Bay LR & California & 1 & $\sim$8 yr & Logistic & PN bloom & --- & No \\
\cite{Lane2009} & & & & regression & (binary) & & \\
\addlinespace
\textbf{DATect} & \textbf{WA/OR} & \textbf{10} & \textbf{21 yr} & \textbf{Ensemble} & \textbf{DA conc.} & \textbf{1 wk} & \textbf{Yes} \\
\textbf{(this study)} & & & & \textbf{(XGB+RF)} & \textbf{+ category} & & \\
\bottomrule
\end{tabular}
\end{table}


% ══════════════════════════════════════════════════════════════════════════════
%  3. METHODOLOGY
% ══════════════════════════════════════════════════════════════════════════════
\section{Methodology}
\label{sec:methodology}

\subsection{Study Area and Data Collection}
\label{sec:data-collection}

\subsubsection{Monitoring Sites}

DATect forecasts DA concentrations at 10 Pacific Coast monitoring sites spanning from Gold Beach, Oregon (42.38\textdegree\,N) to Kalaloch, Washington (47.59\textdegree\,N), a distance of approximately \SI{600}{\kilo\meter} along the coastline (Figure~\ref{fig:study-area}). Sites were selected for their comprehensive historical datasets, consistent sampling over a multi-year period, and spatial representation across the Washington--Oregon coast. Pacific razor clams (\textit{Siliqua patula}) are sampled at all sites. As sessile filter feeders with slow DA depuration (half-life of weeks to months), razor clams continuously accumulate toxin during bloom events, providing an integrated measure of recent DA exposure at each site \cite{Drum1993, Wekell1994, McCabe2023}.

\begin{figure}[H]
\centering
\includegraphics[width=0.55\textwidth]{figures/fig1_study_area.png}
\caption{Study area showing the 10 DATect monitoring sites along the Pacific Northwest coast of Washington and Oregon. Sites span approximately 5\textdegree\ of latitude ($\sim$\SI{600}{\kilo\meter}). The Juan de Fuca eddy region and Heceta Bank, two key bloom retention zones, are indicated.}
\label{fig:study-area}
\end{figure}

Table~\ref{tab:sites} lists the 10 monitoring sites with their coordinates and the number of validation samples available at each site.

\begin{table}[H]
\caption{DATect monitoring sites, geographic coordinates, and validation sample sizes.}
\label{tab:sites}
\small
\begin{tabular}{llccc}
\toprule
\textbf{Site} & \textbf{State} & \textbf{Latitude (\textdegree N)} & \textbf{Longitude (\textdegree W)} & \textbf{$N$ (validation)} \\
\midrule
Kalaloch       & WA & 47.586 & 124.379 & 235 \\
Quinault       & WA & 47.284 & 124.236 & 181 \\
Copalis        & WA & 47.106 & 124.181 & 277 \\
Twin Harbors   & WA & 46.792 & 124.100 & 225 \\
Long Beach     & WA & 46.558 & 124.061 & 209 \\
Clatsop Beach  & OR & 46.029 & 123.917 & 354 \\
Cannon Beach   & OR & 45.882 & 123.959 & 116 \\
Newport        & OR & 44.600 & 124.050 & 218 \\
Coos Bay       & OR & 43.376 & 124.237 & 109 \\
Gold Beach     & OR & 42.377 & 124.414 & 257 \\
\midrule
\multicolumn{4}{l}{\textbf{Total validation samples}} & \textbf{2,181} \\
\bottomrule
\end{tabular}
\end{table}

\subsubsection{Data Sources}

DATect integrates seven categories of environmental and biological data spanning 21 years (2003--2023), summarized in Table~\ref{tab:data-sources}. All data are aligned to a common weekly temporal resolution using Monday-anchored ISO weeks. The dataset is constructed by an automated pipeline that downloads, processes, and integrates all data sources into a unified analysis-ready file.

\begin{table}[H]
\caption{Data sources used in DATect, with temporal resolution, period of record, and operational availability delays enforced during forecasting.}
\label{tab:data-sources}
\small
\begin{tabular}{p{3cm}p{3cm}lll}
\toprule
\textbf{Source} & \textbf{Variables} & \textbf{Resolution} & \textbf{Period} & \textbf{Delay} \\
\midrule
MODIS-Aqua & Chl-$a$, SST, PAR, FLH, K490 & 8-day, $\sim$\SI{4}{\kilo\meter} & 2003--2023 & 7 days \\
(NOAA CoastWatch) & & & & \\
\addlinespace
MODIS anomalies & Chl-$a$ anomaly, SST anomaly & Monthly, $\sim$\SI{4}{\kilo\meter} & 2003--2023 & 2 months \\
\addlinespace
NOAA ERDDAP & PDO, ONI & Monthly & 2003--2023 & 2 months \\
\addlinespace
NOAA ERDDAP & BEUTI & Daily & 2003--2023 & 7 days \\
\addlinespace
USGS NWIS & Columbia River discharge & Daily & 2003--2023 & 7 days \\
\addlinespace
WDOH / ODFW & DA concentrations & Irregular ($\sim$weekly) & 2003--2023 & --- \\
\addlinespace
WDOH / ODFW & \textit{Pseudo-nitzschia} cell counts & Irregular & 2003--2023 & --- \\
\bottomrule
\end{tabular}
\end{table}

\noindent\textit{Note:} Several satellite products in Table~\ref{tab:data-sources} are ingested into the unified dataset for ablation and audit purposes but dropped before final model fitting. In the final training pipeline, chlorophyll-$a$, PAR, K490, chlorophyll anomaly, and site coordinates are removed after confirming negligible marginal value (Appendix~\ref{app:feature-engineering}).

Satellite data from MODIS-Aqua \cite{NASA2002} are retrieved as 8-day composites from the NOAA CoastWatch ERDDAP server for each site's local bounding box. The use of 8-day composites mitigates frequent cloud cover in the PNW. A 7-day processing buffer is enforced to simulate operational satellite data availability delays.

Climate indices include the Pacific Decadal Oscillation (PDO) \cite{Mantua1997}, the Oceanic Ni\~no Index (ONI) \cite{Wolter2011}, and the Biologically Effective Upwelling Transport Index (BEUTI) \cite{Jacox2018}. BEUTI is provided at discrete 1\textdegree\ latitude bands; site-specific values are obtained via inverse distance weighting (IDW) to each site's exact latitude. PDO and ONI are subject to a 2-month reporting delay; BEUTI uses a 7-day buffer. Columbia River discharge from the USGS gauge at The Dalles proxies freshwater input to coastal salinity and nutrient dynamics, with the strongest signal at northern Washington sites.

DA measurements and \textit{Pseudo-nitzschia} cell counts were obtained through public records requests to WDOH, WDFW, ODFW, and the Olympic Region Harmful Algal Bloom (ORHAB) Partnership. Measurements reported as below the detection limit ($< 1$~\si{\microgrampeg}) are imputed as $0.5$~\si{\microgrampeg} (half the detection threshold), a standard practice for left-censored environmental data \cite{Helsel2012}. These biological measurements are irregularly spaced, reflecting operational sampling schedules that intensify during suspected bloom events and relax during quiescent periods.

\subsection{Data Engineering}
\label{sec:data-engineering}

\subsubsection{Weekly Aggregation and Gap Handling}

All data are aggregated to Monday-aligned ISO weeks. DA concentrations within each week are aggregated using the \textit{maximum} operation---a conservative choice that preserves worst-case readings for public health relevance. When multiple measurements exist within a single week (common during active bloom monitoring), the maximum ensures that the highest detected toxin level is retained.

Gaps in the DA and \textit{Pseudo-nitzschia} time series are filled using \textit{exponential decay} interpolation. The key reason to prefer this over linear interpolation is temporal integrity: linear interpolation between two measurements requires the \textit{future} measurement to compute intermediate values, introducing data leakage into the gap-filled training targets. Exponential decay propagates information only forward from the last known observation:
\begin{equation}
\hat{y}(t + \Delta t) = y(t) \cdot e^{-\lambda \Delta t}
\label{eq:decay}
\end{equation}
where $y(t)$ is the last observed value, $\Delta t$ is the gap duration in weeks, and $\lambda$ is the species-specific decay rate ($\lambda_{\text{DA}} = 0.2$/week, half-life $\approx 3.5$ weeks, max gap 2 weeks; $\lambda_{\text{PN}} = 0.3$/week, max gap 4 weeks). This also reflects the biological reality of DA depuration in shellfish tissue \cite{Drum1993, Wekell1994, Bricelj1998}. Gaps exceeding these thresholds are filled with zero, assuming non-detection.

\textbf{Critical distinction: augmented training, raw-only evaluation.} Exponential decay interpolation serves two purposes: (1)~constructing the \textit{environmental feature grid}---the weekly time series from which satellite, climate, and \textit{Pseudo-nitzschia} features are drawn, and (2)~\textit{augmenting the training set} with gap-filled DA values. Training samples include both raw measurements and gap-filled DA values, providing approximately 5$\times$ more training data per site than raw-only training. Where a real measurement exists, it takes priority over the gap-filled value; gap-filled values are used only for weeks without a real sample. This augmentation is controlled by the \texttt{USE\_INTERPOLATED\_TRAINING} configuration flag.

Critically, \textit{test} samples remain restricted to weeks with actual raw DA measurements, ensuring all evaluation metrics reflect predictions of real toxin concentrations. Because the gap-filled values use forward-only exponential decay, they do not introduce temporal leakage into the training set. However, persistence features (\texttt{last\_observed\_da\_raw}, \texttt{weeks\_since\_last\_spike}) and the na\"ive persistence baseline are computed exclusively from real observations. These features should reflect what a coastal manager would actually \textit{know} from field sampling, not synthetic estimates. A \texttt{\_is\_interpolated} marker column tracks which training rows use gap-filled targets, enabling this selective filtering. This design provides the statistical benefits of a larger training set while ensuring that persistence signals and evaluation metrics reflect actual field measurements.

\subsubsection{Derived Environmental Features}

Six derived features encoding domain knowledge---including a marine heatwave flag, BEUTI nonlinear terms, PDO--ONI phase coherence, \textit{Pseudo-nitzschia} transforms, fluorescence efficiency, and a K490 quadratic term---were explored during development (Appendix~\ref{app:feature-engineering}). Individual feature ablation (Table~\ref{tab:feature-ablation}) confirmed that all produce negligible impact ($|\Delta R^2| \leq 0.005$), and all have been removed from the final pipeline.

\subsubsection{Persistence Features, Lag Features, and Rolling Statistics}

DATect also uses \textit{observation-order lag features}---the $N$th most recent raw observation strictly before the test date---rather than fixed time-grid shifts (Appendix~\ref{app:feature-engineering}). Ablation shows these are largely redundant with the persistence features at the one-week horizon ($\Delta R^2 = -0.001$; Section~\ref{sec:results-ablation}).

Two persistence features are computed from raw DA observations:
\begin{itemize}[nosep]
    \item \texttt{last\_observed\_da\_raw}: Forward-fill of the most recent DA measurement
    \item \texttt{weeks\_since\_last\_spike}: Time since last DA $> \SI{20}{\microgrampeg}$ event
\end{itemize}

Rolling statistics (mean, standard deviation, and maximum) are computed over windows of 4, 8, and 12 weeks for each site. Rolling means for 8- and 12-week windows and all rolling standard deviation features have negligible importance and are dropped; only the 4-week rolling mean/maximum and 8-/12-week maximum carry meaningful signal. To prevent the current week's DA from leaking into its own features, rolling statistics are computed on a \texttt{shift(1)} version of the persistence series---that is, the rolling window begins from the \textit{prior} week's forward-filled DA value, not the current week's. This one-row offset ensures that a test point's rolling features reflect only DA history available \textit{before} the prediction date, even if the test point itself corresponds to a raw measurement week.

\subsubsection{Temporal Features and Preprocessing}

Cyclical temporal features (sine/cosine encodings of day-of-year and month) capture seasonal bloom patterns without year-boundary discontinuities. Additional temporal encodings (week-of-year, bloom season flag, linear trend) were explored but confirmed negligible and removed (Appendix~\ref{app:feature-engineering}). All numeric features undergo median imputation for missing values followed by min-max scaling to $[0, 1]$ using scikit-learn \cite{Pedregosa2011}.

\subsection{Model Architecture}
\label{sec:model-selection}

\subsubsection{Two-Model ML Ensemble}

DATect employs a per-site model selection framework, formalized as a weighted two-model ensemble:
\begin{equation}
\hat{y} = w_{\text{XGB}} \cdot \hat{y}_{\text{XGB}} + w_{\text{RF}} \cdot \hat{y}_{\text{RF}}
\label{eq:ensemble}
\end{equation}
where $w_{\text{XGB}} + w_{\text{RF}} = 1$ and the weights are site-specific (Table~\ref{tab:ensemble-config} in Appendix~\ref{app:model-config}). In the current configuration, all sites converged to single-model assignments ($w \in \{0, 1\}$) based on development-set performance, making this effectively a per-site model selection between XGBoost and Random Forest. Blended weights (e.g., 50/50 or 60/40) were tested empirically but produced no improvement over single-model selection on either the held-out or temporal holdout evaluations ($\Delta R^2 \leq -0.004$); the ensemble formulation is retained for generality. Na\"ive persistence is computed for every prediction and reported as a standalone external baseline (Section~\ref{sec:results-overall}), enabling a direct assessment of how much skill the ML models add beyond simple autocorrelation.

\textbf{XGBoost} \cite{Chen2016}: A gradient boosted tree model using the histogram-based training method (\texttt{tree\_method=hist}) with global base hyperparameters ($n_{\text{estimators}} = 400$, max depth $= 6$, learning rate $= 0.05$, subsample $= 0.85$, colsample\_bytree $= 0.85$, colsample\_bylevel $= 0.8$, $\alpha = 0.1$, $\lambda = 1.0$, $\gamma = 0.1$, min\_child\_weight $= 3$) and per-site overrides (Appendix~\ref{app:hyperparams}). During retrospective evaluation, per-anchor hyperparameter selection uses a calibration fraction of 30\% of training data (maximum 20 rows), searching over a site-specific candidate grid of 1--3 configurations. In real-time deployment, per-site hyperparameters are used directly without per-anchor tuning; this simplification is justified because the perturbation study (Appendix~\ref{app:perturbation-sensitivity}) shows that XGBoost performance is insensitive to small hyperparameter variations within the tested range. XGBoost's native handling of missing values and its L1/L2 regularization make it well-suited for the noisy, gap-filled environmental feature space.

\textbf{Random Forest} \cite{Breiman2001}: An ensemble of decision trees with global parameters ($n_{\text{estimators}} = 400$, max depth $= 12$, min samples split $= 5$, min samples leaf $= 3$, max features $= 0.85$). A conservative variant (max depth $= 6$, min samples split $= 10$, max features $= 0.5$) is used at sites with historically low skill. Random Forest's bagging approach provides complementary diversity to XGBoost's sequential boosting.

\textbf{Na\"ive persistence baseline}: Predicts the most recent raw DA measurement at or before the anchor date. Na\"ive persistence is treated as an \textit{external} baseline---it is computed alongside the ensemble for every prediction but excluded from the ensemble blend, providing a clean separation between ML skill and raw autocorrelation. Given the slow depuration kinetics of DA in razor clams \cite{Drum1993}, this baseline is non-trivial and represents a genuinely competitive approach at sites with strong temporal autocorrelation.

\subsubsection{Per-Site Customization}

A distinguishing feature of DATect is that each of the 10 monitoring sites receives a custom configuration encompassing: (1) XGBoost and Random Forest hyperparameter overrides, (2) a curated feature subset, (3) ensemble weights $(w_{\text{XGB}}, w_{\text{RF}})$ in Equation~\ref{eq:ensemble}, and (4) prediction clipping parameters. This per-site approach acknowledges that the 10 sites exhibit fundamentally different DA dynamics due to their distinct oceanographic settings, sampling histories, and bloom characteristics. A single global model would be forced to compromise across these diverse regimes, while per-site customization allows each site's model to specialize.

Table~\ref{tab:ensemble-config} in Appendix~\ref{app:model-config} summarizes the per-site configuration. In the final paper configuration, six sites select Random Forest and four select XGBoost, so the ``ensemble'' acts as a site-wise model selector rather than a blended average. Blended weights were tested but produced slightly worse results than single-model selection ($\Delta R^2 = -0.004$ standard, $-0.010$ temporal holdout), consistent with the two models being too correlated on the same feature set to benefit from averaging. Most moderate- and low-skill sites also use stronger clipping or regularization, reflecting the need for more conservative predictions when data are sparse or highly variable. Na\"ive persistence remains available as a standalone model type for direct baseline comparison (Section~\ref{sec:results-overall}).

\textbf{Configuration selection procedure.} Per-site ensemble weights, feature subsets, and hyperparameter overrides were selected through an iterative manual process informed by both domain knowledge and preliminary evaluation results. For each site, XGBoost and Random Forest were evaluated independently using the leak-free retrospective protocol described in Section~\ref{sec:validation}, using a \textit{development} test set (seed = 42, 20\% sampling fraction). Ensemble weights were set proportionally to each component's relative per-site performance; sites where one model clearly dominated converged to a single-model configuration. Blended weights were also evaluated but did not improve over single-model selection (Appendix~\ref{app:weight-robustness}). Na\"ive persistence was also evaluated independently as an external baseline but was excluded from the ensemble blend in order to provide a clean separation between ML ensemble skill and raw autocorrelation. Feature subsets were curated by removing feature groups that degraded site-specific performance, prioritizing parsimony at data-sparse sites to reduce overfitting.

\subsubsection{Competitor Models}

Two additional models are included as competitors (not baselines):

\textbf{Ridge regression} \cite{Hoerl1970}: L2-regularized linear regression ($\alpha = 1.0$) using the full feature set, representing the best linear model.

\textbf{Logistic regression}: For classification tasks, using an L-BFGS solver with $C = 1.0$ and 1,000 maximum iterations.

\subsubsection{Classification}
\label{sec:classification}

DATect classifies continuous DA predictions into four risk categories aligned with fixed regulatory thresholds: Low (0--5), Moderate (5--20), High (20--40), and Extreme ($>$40) \si{\microgrampeg} (Appendix~\ref{app:classification-details}, Table~\ref{tab:da-categories}). The 20~\si{\microgrampeg} threshold between Moderate and High corresponds to the FDA regulatory action level---the concentration above which shellfish harvest is prohibited. All classification results reported in this paper use this threshold-based approach applied to regression predictions. The web dashboard additionally trains a dedicated XGBoost classifier to provide per-class probability estimates for visualization purposes.

\subsubsection{Confidence Intervals}

Prediction uncertainty is summarized via \textbf{quantile regression} \cite{Koenker1978}, which trains separate XGBoost models with the \texttt{reg:quantile} objective at specified percentiles---the 5th, 50th, and 95th for real-time forecasts (90\% interval) and the 10th, 50th, and 90th for retrospective summaries (80\% interval). In ensemble mode, these XGBoost quantiles are combined with point-estimate Random Forest and na\"ive components using the site-specific ensemble weights, producing an empirical predictive band rather than a fully calibrated posterior interval. Bootstrap summaries reported elsewhere in the paper similarly quantify variation across held-out prediction rows, but they do not eliminate all temporal dependence in the retrospective frame.

\subsection{Validation and Scientific Integrity}
\label{sec:validation}

\subsubsection{Leak-Free Validation Protocol}

DATect enforces nine explicit data-leakage prevention mechanisms, summarized in Table~\ref{tab:leakage-safeguards}. These address the pervasive problem of temporal data leakage in environmental ML \cite{Kapoor2023, Roberts2017}. Unlike many published environmental ML studies that rely on random train-test splits, DATect enforces strict chronological ordering at every stage of the pipeline.

\begin{table}[H]
\caption{Nine data-leakage prevention mechanisms enforced in DATect. Each safeguard is structural (enforced in code) rather than procedural.}
\label{tab:leakage-safeguards}
\small
\begin{tabular}{clp{7cm}}
\toprule
\textbf{\#} & \textbf{Safeguard} & \textbf{Implementation} \\
\midrule
1 & Chronological split & Training data restricted to date $\leq$ anchor date \\
2 & Forward-only gap-filling & Exponential decay propagates only from past observations; unlike linear interpolation, gap-filled values never depend on future measurements \\
3 & Observation-order lags & Past-only shifts on raw DA measurements \\
4 & Satellite processing delay & 7-day buffer enforced at data ingestion \\
5 & Climate reporting delay & 2-month buffer for PDO, ONI, monthly anomalies \\
6 & Fixed category thresholds & DA risk bins (5, 20, 40 \si{\microgrampeg}) are regulatory constants, not derived from training data \\
7 & Persistence recomputation & Persistence features recomputed per test point using only training data \\
8 & Fresh model per test point & New model trained for each prediction; no shared state \\
9 & Runtime assertion & \texttt{\_verify\_no\_data\_leakage()} called on every prediction \\
\bottomrule
\end{tabular}
\end{table}

The central concept is the \textit{anchor date}, defined as $t_{\text{anchor}} = t_{\text{forecast}} - 7$ days. All training data must satisfy date $\leq t_{\text{anchor}}$, and the test date must satisfy $t_{\text{test}} > t_{\text{anchor}}$. This is verified by an assertion function called on every single prediction, raising an error if violated. Safeguard \#8 (fresh model per test point) is particularly aggressive---rather than training a single model and applying it to multiple test points, DATect retrains from scratch for each prediction, ensuring that no information from future test points can influence earlier predictions through model state.

\subsubsection{Retrospective Evaluation}

Model performance is evaluated through leak-free retrospective validation using a two-stage protocol. Per-site configurations (ensemble weights, feature subsets, hyperparameter overrides) were developed using a \textit{development} test set (seed = 42, 20\% sampling fraction). All results reported in this paper are evaluated on a separate \textit{held-out} test set drawn with a different random seed (seed = 123) and a larger sampling fraction (40\%), ensuring that no reported metric informed configuration selection. For each site, 40\% of raw DA measurements are sampled as candidate test points, subject to a history requirement that $\geq 33\%$ of the site's total record must precede the test date. At each test point, a fresh model is trained on all data up to the anchor date, predictions are generated, and metrics are computed. Bootstrap 95\% confidence intervals ($B = 10{,}000$; percentile method \cite{Efron1979}) are reported for key metrics as summaries of variation across held-out prediction rows, particularly at sites with limited sample sizes.

As an additional validation, we report a \textit{temporal holdout} evaluation in which test points are restricted to 2019--2023 ($N = 1{,}036$, Table~\ref{tab:temporal-holdout}). Per-site configurations were developed on the seed-42 development set, which samples 20\% of raw measurements from all years; by the sampling design, approximately 20\% of the 1,036 temporal holdout test points ($\sim$207 observations) also appeared as development test points, while the remaining $\sim$829 (80\%) are encountered here for the first time. A five-seed perturbation study (Appendix~\ref{app:perturbation-sensitivity}) confirms that all tuning choices are robust across seeds, mitigating any residual influence from the $\sim$20\% overlap. Training for each test point still uses all data before the anchor date, so the hold-out designation affects only which dates appear as test points, not the expanding-window training mechanism.

Algorithm~\ref{alg:forecast} presents the pseudocode for a single forecast generation, and Figure~\ref{fig:ml-pipeline} provides a visual summary of the same chronology-preserving workflow.

\begin{algorithm}[H]
\SetAlgoLined
\KwIn{forecast\_date $t$, site $s$, model\_type}
\KwOut{Prediction $\hat{y}$, confidence interval, risk category}
$t_{\text{anchor}} \gets t - 7$ days\;
Load feature frame for site $s$\;
$\mathcal{D}_{\text{train}} \gets \{(x_i, y_i) : \text{date}_i \leq t_{\text{anchor}}\}$\;
$x_{\text{test}} \gets$ environmental features at or before $t_{\text{anchor}}$\;
Recompute persistence features from $\mathcal{D}_{\text{train}}$ only\;
Build observation-order lag features from $\mathcal{D}_{\text{train}}$ only\;
Apply preprocessing (median impute $\to$ MinMax scale)\;
\texttt{\_verify\_no\_data\_leakage}($\mathcal{D}_{\text{train}}, t, t_{\text{anchor}}$)\;
$m \gets$ site-specific model type for $s$ \tcp*{XGBoost or Random Forest}
Train model $m$ on $\mathcal{D}_{\text{train}}$ with site-specific hyperparameters\;
$\hat{y} \gets m(x_{\text{test}})$\;
$\hat{y}_{\text{Naive}} \gets$ most recent DA in $\mathcal{D}_{\text{train}}$ \tcp*{external baseline only}
Apply prediction clipping (site-specific quantile and/or hard max)\;
Compute confidence interval via quantile regression\;
Classify $\hat{y}$ into risk category (Low/Moderate/High/Extreme)\;
\Return{$\hat{y}$, CI, category, feature importance}\;
\caption{DATect Single Forecast Generation}
\label{alg:forecast}
\end{algorithm}

\begin{figure}[H]
\centering
\includegraphics[width=0.85\textwidth]{figures/fig3b_ml_pipeline.png}
\caption{DATect forecasting pipeline. For each forecast point, an expanding window enforces chronological splitting at the anchor date and trains a fresh site-specific XGBoost or Random Forest model. Features include observation-order lags, rolling statistics, persistence indicators recomputed from real observations only, and environmental covariates available at or before the anchor date. A separate spike classifier supports transition-focused early warning, and post-processing applies prediction clipping, empirical predictive bands, and threshold-based DA risk classification.}
\label{fig:ml-pipeline}
\end{figure}

\subsubsection{Evaluation Metrics}

Regression performance is assessed using the coefficient of determination ($R^2$), mean absolute error (MAE), and root mean squared error (RMSE). Classification performance uses accuracy, per-class precision, recall, and F1 score. Spike detection (binary DA $> \SI{20}{\microgrampeg}$) is evaluated using the F2 score, which weights recall twice as heavily as precision ($\beta = 2$), and a recall-optimized operating point ($p \geq 0.10$), reflecting the asymmetric cost structure where a missed toxic event (potential human illness) carries far greater consequence than a false alarm (unnecessary precautionary advisory).

DATect is deployed as a publicly accessible web dashboard providing real-time forecasting and retrospective analysis capabilities. Full architectural details, deployment configuration, and visualization descriptions are provided in Appendix~\ref{app:webapp}.


% ══════════════════════════════════════════════════════════════════════════════
%  4. RESULTS
% ══════════════════════════════════════════════════════════════════════════════
\section{Results}
\label{sec:results}

\subsection{Overall Model Performance}
\label{sec:results-overall}

Table~\ref{tab:model-comparison} compares the five model variants evaluated under leak-free retrospective validation on the held-out test set ($n = 2{,}181$). The two-model ML ensemble achieves $R^2 = 0.215$, MAE $= \SI{6.42}{\microgrampeg}$, and RMSE $= \SI{17.79}{\microgrampeg}$. Figure~\ref{fig:site-performance} visualizes the site-to-site split that drives this pooled result.

\begin{table}[H]
\caption{Overall regression performance of five model variants under leak-free retrospective validation on the held-out test set ($n = 2{,}181$, seed = 123). Bootstrap 95\% confidence intervals ($B = 10{,}000$; percentile method) are shown for $R^2$ and MAE and summarize variation across held-out prediction rows. Paired $\Delta R^2$ vs.\ ensemble reported with 95\% CI.}
\label{tab:model-comparison}
\small
\begin{tabular}{lcccc}
\toprule
\textbf{Model} & $\boldsymbol{R^2}$ \textbf{[95\% CI]} & \textbf{MAE [95\% CI] (\si{\microgrampeg})} & \textbf{RMSE (\si{\microgrampeg})} & \textbf{$\Delta R^2$ vs.\ ens.\ [95\% CI]} \\
\midrule
\textbf{Ensemble (XGB+RF)} & \textbf{0.215} [0.096, 0.342] & \textbf{6.42} [5.78, 7.15] & \textbf{17.79} & --- \\
XGBoost                    & 0.221 [0.112, 0.339] & 6.38 [5.67, 7.08] & 17.72 & $-$0.006 [$-$0.020, $+$0.007] \\
Random Forest              & 0.190 [0.057, 0.320] & 6.47 [5.79, 7.25] & 18.07 & $+$0.025 [$+$0.005, $+$0.045] \\
Linear (Ridge)             & 0.202 [0.070, 0.324] & 6.62 [5.93, 7.39] & 17.94 & $+$0.013 [$-$0.020, $+$0.042] \\
Na\"ive Persistence        & $-$0.426 [$-$0.951, $-$0.034] & 7.73 [6.84, 8.76] & 23.98 & $+$0.641 [$+$0.297, $+$1.122] \\
\bottomrule
\end{tabular}
\end{table}

XGBoost achieves the highest individual model $R^2 = 0.221$, marginally above the ensemble ($R^2 = 0.215$). That the ensemble slightly underperforms XGB-everywhere reflects that per-site RF selection at several WA sites trades a small amount of pooled $R^2$ for better per-site calibration; this is a consequence of pooled $R^2$ not being a weighted average of per-site values. Random Forest ($R^2 = 0.190$) and Linear/Ridge ($R^2 = 0.202$) trail XGB but remain competitive. Na\"ive persistence achieves $R^2 = -0.426$ at the pooled held-out scale, indicating that simply predicting the most recent DA observation is not sufficient once all sites are combined. This pooled result does not imply that persistence is useless at every site; rather, it underscores the importance of per-site modeling under strong regional heterogeneity. Paired bootstrap differences between the ensemble and individual ML models are small and statistically consistent with zero (Table~\ref{tab:model-comparison}), while the ensemble's advantage over na\"ive persistence is large and clearly significant ($\Delta R^2 = +0.641$\, $[+0.297,\, +1.122]$), reflecting the value of per-site model selection and feature engineering over raw autocorrelation.

\subsection{Per-Site Regression Performance}
\label{sec:results-per-site}

Table~\ref{tab:per-site-results} presents detailed per-site ensemble performance. Sites naturally partition into three tiers:

\begin{table}[H]
\caption{Per-site ensemble regression and classification performance on the held-out test set (seed = 123, separated from configuration-tuning data). Bootstrap 95\% confidence intervals ($B = 10{,}000$) are shown for $R^2$ and MAE. Sites are ordered by $R^2$ and grouped into performance tiers. Total test set $n = 2{,}181$.}
\label{tab:per-site-results}
\small
\begin{tabular}{lrrllrl}
\toprule
\textbf{Site} & $\boldsymbol{N}$ & $\boldsymbol{R^2}$ \textbf{[95\% CI]} & \textbf{MAE [95\% CI]} & \textbf{RMSE} & \textbf{Acc.} & \textbf{Tier} \\
& & & \textbf{(\si{\microgrampeg})} & \textbf{(\si{\microgrampeg})} & & \\
\midrule
Copalis        & 277 & \textbf{0.789} [0.710, 0.844] & 2.81 [2.32, 3.40] & 5.28 & 0.823 & High-skill \\
Twin Harbors   & 225 & \textbf{0.594} [0.473, 0.764] & 4.66 [3.36, 6.21] & 12.20 & 0.809 & High-skill \\
Quinault       & 181 & \textbf{0.582} [0.498, 0.736] & 4.13 [2.96, 5.58] & 10.15 & 0.762 & High-skill \\
Long Beach     & 209 & \textbf{0.632} [0.560, 0.710] & 4.47 [3.51, 5.51] & 8.63 & 0.746 & High-skill \\
Kalaloch       & 235 & \textbf{0.480} [0.349, 0.641] & 3.78 [2.73, 5.09] & 10.14 & 0.813 & High-skill \\
\addlinespace
Clatsop Beach  & 354 & 0.292 [$-$0.229, 0.562] & 6.14 [4.85, 7.55] & 14.04 & 0.706 & Moderate \\
Gold Beach     & 257 & 0.051 [$-$0.113, 0.236] & 8.02 [5.30, 10.78] & 23.45 & 0.763 & Moderate \\
\addlinespace
Cannon Beach   & 116 & $-$0.045 [$-$0.091, $-$0.010] & 3.93 [1.38, 6.97] & 16.15 & 0.922 & Low-skill \\
Coos Bay       & 109 & $-$0.038 [$-$0.681, 0.304] & 24.55 [18.75, 30.47] & 40.12 & 0.450 & Low-skill \\
Newport        & 218 & $-$0.243 [$-$1.512, $-$0.017] & 10.30 [7.28, 14.13] & 28.28 & 0.610 & Low-skill \\
\midrule
\textbf{Overall} & \textbf{2,181} & \textbf{0.215} [0.096, 0.342] & \textbf{6.42} [5.78, 7.15] & \textbf{17.79} & \textbf{0.751} & --- \\
\bottomrule
\end{tabular}
\vspace{0.3em}\\
\footnotesize{Adjusted $R^2$ values (not shown) differ from $R^2$ by $< 0.02$ at all sites; the modest sample-size-to-feature ratios at data-sparse sites are mitigated by per-site feature subsetting (Appendix~\ref{app:features}).}
\end{table}

\begin{table}[H]
\caption{Per-site ensemble performance on the 2019--2023 temporal holdout ($N = 1{,}036$). Per-site configurations were developed on the seed-42 development set, which draws test points from all years; however, the development set's 20\% sampling means that the large majority of post-2019 observations appear here for the first time, and a five-seed perturbation study (Appendix~\ref{app:perturbation-sensitivity}) confirms that all tuning choices are stable across seeds. Training for each test point uses all data before the anchor date. Pooled $R^2 = 0.315$ [0.138, 0.476], WA mean $R^2 = 0.549$, OR mean $R^2 = 0.113$.}
\label{tab:temporal-holdout}
\small
\begin{tabular}{lrllrl}
\toprule
\textbf{Site} & $\boldsymbol{N}$ & $\boldsymbol{R^2}$ \textbf{[95\% CI]} & \textbf{MAE [95\% CI] (\si{\microgrampeg})} & \textbf{RMSE (\si{\microgrampeg})} & \textbf{Tier} \\
\midrule
Quinault      & 144 & \textbf{0.735} [0.609, 0.831] & 2.72 [2.12, 3.40] & 4.69  & High-skill \\
Copalis       & 152 & \textbf{0.629} [0.504, 0.753] & 3.20 [2.51, 3.96] & 5.56  & High-skill \\
Long Beach    & 101 & \textbf{0.628} [0.452, 0.766] & 5.27 [3.83, 6.83] & 9.33  & High-skill \\
Twin Harbors  & 103 & \textbf{0.461} [0.236, 0.712] & 5.18 [3.67, 7.04] & 10.26 & High-skill \\
Clatsop Beach & 118 & \textbf{0.522} [0.013, 0.770] & 9.78 [7.13, 12.73] & 18.69 & High-skill$^\dagger$ \\
Coos Bay      &  89 & \textbf{0.616} [0.181, 0.857] & 8.08 [5.32, 11.55] & 16.78 & High-skill$^\dagger$ \\
\addlinespace
Kalaloch      &  82 & 0.289 [0.138, 0.612] & 2.55 [1.44, 3.92] & 6.31  & Moderate \\
Cannon Beach  &  13 & 0.000$^\ddagger$ & 0.03 & 0.03  & N/A$^\ddagger$ \\
Newport       & 124 & $-$0.305 [$-$1.514, 0.129] & 12.88 [9.19, 17.16] & 25.15 & Low-skill \\
Gold Beach    & 110 & $-$0.271 [$-$2.336, $-$0.048] & 10.56 [6.95, 15.04] & 24.96 & Low-skill \\
\midrule
\textbf{Overall} & \textbf{1,036} & \textbf{0.315} [0.138, 0.476] & \textbf{6.55} [5.72, 7.43] & \textbf{15.35} & --- \\
\bottomrule
\end{tabular}
\vspace{0.3em}\\
\footnotesize{Bootstrap 95\% CIs ($B = 10{,}000$; percentile method). $^\dagger$Clatsop Beach and Coos Bay shift from Moderate/Low-skill on the standard eval to High-skill on the temporal holdout ($R^2 > 0.4$). Coos Bay's improvement ($R^2 = 0.616$ vs.\ $-0.038$) is particularly striking and may reflect a post-2019 regime shift in DA dynamics consistent with the 2019--2021 marine heatwave. $^\ddagger$Cannon Beach ($N = 13$) has insufficient post-2019 measurements for stable per-site estimates; essentially all 13 holdout observations are near-zero DA, making this a degenerate sample rather than genuine forecast skill. Bootstrap CI omitted.}
\end{table}

\begin{table}[H]
\caption{Persistence-forecast $R^2$ ceiling per site, computed as the squared lag-1 autocorrelation ($\rho^2$) of consecutive raw DA measurements. The ceiling bounds the maximum $R^2$ achievable by a \emph{persistence-only} (last-value) one-step-ahead forecast; models that incorporate environmental predictors (SST, BEUTI, PDO, discharge, etc.) can and do exceed $\rho^2$, as illustrated by Gold Beach ($\rho^2 = 0.002$, held-out $R^2 = 0.051$). Sites with near-zero ceilings indicate that persistence alone carries very little information about future DA levels.}
\label{tab:predictability-ceiling}
\small
\begin{tabular}{lrrrr}
\toprule
\textbf{Site} & $\boldsymbol{N}$ & $\boldsymbol{\rho}$ \textbf{(lag-1)} & $\boldsymbol{R^2}$ \textbf{ceiling} ($\rho^2$) & \textbf{Held-out} $\boldsymbol{R^2}$ \\
\midrule
Copalis       & 834 & 0.909 & 0.826 & 0.789 \\
Kalaloch      & 654 & 0.883 & 0.780 & 0.480 \\
Long Beach    & 698 & 0.881 & 0.776 & 0.632 \\
Twin Harbors  & 690 & 0.863 & 0.744 & 0.594 \\
Quinault      & 564 & 0.836 & 0.698 & 0.582 \\
\addlinespace
Clatsop Beach & 1087 & 0.501 & 0.251 & 0.292 \\
Cannon Beach  &  305 & 0.443 & 0.197 & $-$0.045 \\
Newport       &  708 & 0.191 & 0.036 & $-$0.243 \\
Coos Bay      &  335 & 0.054 & 0.003 & $-$0.038 \\
Gold Beach    &  717 & 0.038 & 0.002 & 0.051 \\
\bottomrule
\end{tabular}
\vspace{0.3em}\\
\footnotesize{$\rho$ computed on consecutively paired raw DA observations (non-interpolated), excluding gaps $> 28$ days. WA mean ceiling: 0.765. OR mean ceiling: 0.098. Held-out $R^2$ from Table~\ref{tab:per-site-results} (seed = 123).}
\end{table}

All five Washington sites fall in the \textbf{high-skill} tier ($R^2 = 0.480$--$0.789$), confirming that strong performance is not an artifact of per-site tuning; Copalis leads at $R^2 = 0.789$ [0.710, 0.844], and Figure~\ref{fig:pred-vs-obs} illustrates how the ensemble tracks observed DA dynamics at this site over the full 21-year record. Two Oregon sites show \textbf{moderate} skill: Clatsop Beach ($R^2 = 0.292$) and Gold Beach ($R^2 = 0.051$) are marginally positive. Cannon Beach, Coos Bay, and Newport form the \textbf{low-skill} tier ($R^2 \leq 0$). Cannon Beach ($R^2 = -0.045$) achieves 92.2\% classification accuracy only because nearly all readings are Low; Newport's wide confidence interval [$-1.512$, $-0.017$] reflects episodic, unpredictable bloom events (Table~\ref{tab:predictability-ceiling}).

The 2019--2023 temporal holdout ($n = 1{,}036$, Table~\ref{tab:temporal-holdout}) provides additional validation. Pooled $R^2 = 0.315$ exceeds the standard result (0.215), as the post-2019 period contains proportionally more Washington test points and stronger autocorrelation. Washington sites achieve mean $R^2 = 0.549$; Oregon sites average $R^2 = 0.113$. The notable anomaly is Coos Bay ($R^2 = 0.616$ vs.\ $-0.038$ on the standard eval), possibly reflecting a post-2019 regime shift during the 2019--2021 marine heatwave.

\begin{figure}[H]
\centering
\includegraphics[width=0.9\textwidth]{figures/fig2_site_performance.png}
\caption{Held-out per-site regression performance. Washington sites (blue) consistently outperform Oregon sites (red), reinforcing the interpretation that monitoring density and site-specific data quality are the main drivers of forecast skill in the current record. Point labels denote the number of validation samples at each site.}
\label{fig:site-performance}
\end{figure}

\begin{figure}[H]
\centering
\includegraphics[width=0.9\textwidth]{figures/fig5_pred_vs_obs.png}
\caption{Ensemble predictions versus observed DA concentrations at Copalis, the highest-skill site. (a)~Time series of observed DA (black dots) and ensemble predictions (blue line) over the 21-year record, with the FDA regulatory threshold at \SI{20}{\microgrampeg} (red dashed line). (b)~Scatter plot of predicted versus observed values with the 1:1 line shown. Evaluated on the development set (seed = 42, $n = 148$ at this site); held-out metrics ($R^2 = 0.789$) are reported in Table~\ref{tab:per-site-results}.}
\label{fig:pred-vs-obs}
\end{figure}

\subsection{Classification Performance}
\label{sec:results-classification}

Threshold-based four-category performance is summarized in Appendix~\ref{app:classification-details} (Tables~\ref{tab:classification} and \ref{tab:confusion}). Overall accuracy is 75.1\% on the held-out test set, driven largely by reliable identification of the Low category (F1 $= 0.862$). Higher-risk categories remain more difficult because of limited samples and class imbalance. Most errors occur between adjacent categories, but 15 of 108 Extreme samples (13.9\%) are misclassified as Low---a safety-relevant failure mode that motivates the dedicated spike analysis below and is revisited in Section~\ref{sec:discussion-limitations}.

\subsection{Spike Transition Detection}
\label{sec:results-spike}

While regression metrics characterize average forecasting skill, the operationally critical question is whether DATect can anticipate \emph{spike transitions}---the moment DA crosses from below to at or above the \SI{20}{\microgrampeg} regulatory threshold. These below-to-above crossings are exactly the events that trigger beach closures, and early detection would provide the advance warning that current reactive monitoring cannot.

\subsubsection{Event Cohorts and Definitions}

Two related but distinct event cohorts are used in this section. Environmental lead-signal analysis uses all raw-observation threshold crossings identified from consecutive site-level DA measurements (385 events), because the goal is to ask whether pre-spike environmental conditions differ from baseline in the underlying record. Forecast evaluation uses only held-out retrospective test rows that correspond to below-to-above-threshold crossings relative to the most recent prior raw observation at the same site (89 events, seed = 123). The first cohort is used for mechanistic evidence; the second is used for operational skill evaluation.

\subsubsection{Environmental Lead Signals}
\label{sec:results-env-lead}

Analysis of 385 spike transition events across all 10 sites reveals that environmental features change detectably in the 1--4 weeks \emph{before} DA crosses the regulatory threshold. The Pacific Decadal Oscillation (PDO) showed the strongest pre-spike signal with rank-biserial effect size $r = 0.50$, followed by SST anomaly ($r = 0.40$), BEUTI ($r = 0.32$), and Columbia River discharge ($r = 0.27$, Washington sites only). All seven environmental features tested showed statistically significant deviations from baseline ($p < 0.05$, Wilcoxon signed-rank test) during the four weeks preceding spike transitions. These lead signals confirm that the environmental information needed for anticipatory spike detection exists in the feature set.

\subsubsection{Regression Threshold Limitations}
\label{sec:results-threshold-limit}

Despite this, the regression ensemble's spike transition recall at the standard \SI{20}{\microgrampeg} decision threshold is poor: only 12.4\% of held-out transition events are correctly predicted as spikes ($n = 89$, seed = 123). Na\"ive persistence achieves 23.6\%. The ensemble \emph{does} predict elevated values before spikes---mean predicted DA for transition events is \SI{13.1}{\microgrampeg}---but the conservative bias inherent in MSE-optimized regression prevents predictions from crossing the regulatory threshold. A threshold sensitivity analysis confirms this: lowering the decision boundary from 20 to \SI{12}{\microgrampeg} increases transition recall substantially while maintaining acceptable precision.

\subsubsection{Transition-Focused Binary Classifier}
\label{sec:results-spike-classifier}

To exploit the pre-spike environmental signals without the regression's conservative bias, we developed a dedicated binary spike classifier trained alongside the regression ensemble. The classifier uses a \emph{safe-baseline} training strategy: it trains only on observations where the previous DA measurement was below \SI{20}{\microgrampeg}, forcing it to learn from environmental and rolling features rather than simple persistence. This prevents the classifier from relying on ``last DA was already high'' as its primary signal---a strategy that would detect ongoing spikes but miss the operationally critical below-to-above transitions.

Table~\ref{tab:spike-transition} compares spike transition detection across approaches on the held-out test set (seed = 123, $n = 2{,}181$, 89 transitions). Figure~\ref{fig:spike-summary} highlights the same operating-point tradeoff visually.

\begin{table}[H]
\caption{Spike transition detection performance on the held-out test set ($n = 2{,}181$, 89 below-to-above-threshold transitions). Transition recall measures the fraction of below-to-above-threshold crossings correctly identified. The hybrid alert (spike classifier combined with regression trigger at \SI{12}{\microgrampeg}) achieves 3.1$\times$ the transition recall of na\"ive persistence.}
\label{tab:spike-transition}
\centering
\begin{tabular}{lcccc}
\toprule
\textbf{Approach} & \textbf{Transition Recall} & \textbf{Spike Recall} & \textbf{Precision} & \textbf{F2} \\
\midrule
Na\"ive persistence              & 0.236 & 0.754 & 0.546 & 0.700 \\
Ensemble ($\geq$20 \si{\microgrampeg}) & 0.124 & 0.558 & 0.618 & 0.569 \\
Hybrid alert (classifier + \SI{12}{\microgrampeg} trigger) & \textbf{0.734} & \textbf{0.859} & 0.344 & 0.661 \\
\bottomrule
\end{tabular}
\vspace{0.3em}\\
\footnotesize{F2 weights recall twice as heavily as precision ($\beta = 2$). The hybrid alert's F2 (0.661) is lower than na\"ive persistence (0.700) because its precision drops sharply; however, the hybrid's advantage is specifically in \textit{transition} recall (0.734 vs.\ 0.236), which F2 does not separately capture. For closure-relevant early warning, transition recall is the operationally critical metric.}
\end{table}

\begin{figure}[H]
\centering
\includegraphics[width=0.88\textwidth]{figures/fig4_spike_detection_summary.png}
\caption{Spike-focused summary of the held-out results in Table~\ref{tab:spike-transition}. Thresholding the regression ensemble at \SI{20}{\microgrampeg} misses most below-to-above-threshold transitions, whereas the safe-baseline classifier substantially improves transition recall at the cost of lower precision. The hybrid alert used in the deployed dashboard is discussed in Section~\ref{sec:discussion-spike}.}
\label{fig:spike-summary}
\end{figure}

The hybrid alert---combining the transition-focused classifier with a regression trigger at \SI{12}{\microgrampeg}---achieves 73.4\% transition recall compared to 23.6\% for na\"ive persistence and 12.4\% for the standard regression ensemble. This 3.1$\times$ improvement in transition recall ($0.734 / 0.236$) comes at the cost of substantially lower precision (0.344 vs.\ 0.546), yielding a lower overall F2 score (0.661 vs.\ 0.700 for na\"ive persistence). The tradeoff is intentional: in the closure-warning context, the cost of a missed transition (delayed public health response) far exceeds the cost of a false alert (increased sampling attention), making transition recall the operationally critical metric rather than F2. The classifier's top features are rolling DA maximum over 4 weeks (gain importance 0.41--0.45), consistent with the biological mechanism where gradual toxin accumulation precedes a threshold crossing (Figure~\ref{fig:feature-importance}). Environmental features including PDO, BEUTI, and fluorescence contribute meaningfully when persistence features are excluded by the safe-baseline filter.

\subsection{Ablation Study}
\label{sec:results-ablation}

Component ablation is summarized in Appendix~\ref{app:component-ablation} (Table~\ref{tab:ablation}) using the development set (seed = 42, $n = 1{,}177$). Per-site customization is the most impactful component ($\Delta R^2 = -0.158$), confirming that the 10 monitoring sites cannot be represented well by a single global configuration. Interpolated training provides a modest improvement ($\Delta R^2 = -0.029$ when removed), while observation-order lags and derived features show negligible standalone impact at the one-week horizon. These latter effects are best interpreted as redundancy rather than as proof that the features are uninformative: persistence and rolling statistics already capture much of the same recent-history signal.

A fine-grained per-feature ablation (Appendix~\ref{app:feature-ablation}, Table~\ref{tab:feature-ablation}) confirms that no individual environmental feature exceeds noise-level impact ($|\Delta R^2| \leq 0.005$), and removing all low-importance features simultaneously costs only $\Delta R^2 = -0.004$. High feature importance does not imply high marginal contribution: lag features rank highly in model splits but are largely redundant with \texttt{last\_observed\_da\_raw} ($\Delta R^2 = -0.001$ when removed). A targeted perturbation study (Appendix~\ref{app:perturbation-sensitivity}, Tables~\ref{tab:perturbation}--\ref{tab:seed-stability}) confirms that RF hyperparameters, clipping thresholds, and per-site model selection are all well-calibrated, with Washington sites stable across five seeds (mean $R^2 = 0.61 \pm 0.05$) and Oregon variance reflecting data scarcity rather than tuning fragility.


\subsection{Diagnostic Findings}
\label{sec:results-temporal-holdout}

Two diagnostic analyses characterize where the ensemble fails. An episodic-error decomposition shows that 88.6\% of total squared error concentrates in the top 5\% of worst-performing test weeks ($n = 110$ rows)---all Oregon observations exceeding \SI{100}{\microgrampeg}, consistent with near-zero autocorrelation ceilings at those sites (Table~\ref{tab:predictability-ceiling}). A bias analysis in the 12--30~\si{\microgrampeg} transition zone ($n = 311$ observations) reveals a systematic $+2.2$~\si{\microgrampeg} under-prediction, a predictable consequence of MSE-optimized regression on a right-skewed target and the primary motivation for the hybrid alert's transition-focused classifier.

% ══════════════════════════════════════════════════════════════════════════════
%  5. DISCUSSION
% ══════════════════════════════════════════════════════════════════════════════
\section{Discussion}
\label{sec:discussion}

\subsection{Data Quality and Site Heterogeneity}

The clearest pattern in the held-out results is the Washington--Oregon performance gradient. All five Washington sites achieve $R^2 = 0.480$--$0.789$, whereas Oregon sites range from modest positive skill to strongly negative $R^2$. The most parsimonious explanation is data quality: Washington records from WDFW and Tribal monitoring are denser and more temporally regular, so persistence, lag, and rolling features can be estimated more reliably. Oregon records are sparser and less regular, reducing the effective information content available to both linear and nonlinear models.

Genuine oceanographic differences may also contribute: Washington beaches are influenced by the Juan de Fuca eddy and the Columbia River plume, both of which may increase DA persistence or provide stronger environmental structure \cite{Trainer2009, McCabe2023}. Even so, monitoring density is currently the binding constraint on forecast skill, and better-sampled Oregon records would likely improve performance more than a more complex model family.

\subsection{Persistence, Model Choice, and ML Contribution}

The difference between development-set $R^2$ (0.414, seed = 42) and held-out test-set $R^2$ (0.215, seed = 123) merits interpretation. First, the gap is expected by design: per-site configurations were tuned on seed = 42, so reporting on a separate seed produces less biased estimates. Second, the held-out set's larger sampling fraction (40\% vs.\ 20\%) includes more Oregon test points, and the four low-skill Oregon sites disproportionately depress the aggregate. Washington sites achieve $R^2 = 0.480$--$0.789$ on both seeds, confirming that tuned configurations generalize where sufficient data quality exists; the aggregate $R^2$ reflects Oregon data sparsity, not overfitting.

Strong persistence at several Washington sites is biologically plausible. DA depuration from razor clam tissue occurs over weeks to months \cite{Drum1993, Bricelj1998}, so one-week-ahead DA often remains close to the most recent observation during sustained bloom periods. DATect therefore treats na\"ive persistence as a serious external baseline rather than a straw man. The pooled held-out result ($R^2 = -0.426$) should not be read as evidence that persistence is never useful; rather, it reflects the fact that persistence is highly site-dependent and degrades sharply at the low-skill Oregon sites.

The near-tie between the tree-based models and Ridge regression further shows that most forecast signal comes from the feature set and the per-site design, not from an intrinsically superior model family. Removing per-site customization costs $\Delta R^2 = -0.158$ in the development ablation, far more than any other single component. Likewise, model-selection robustness in Appendix~\ref{app:weight-robustness} shows that choosing XGBoost versus Random Forest per site changes aggregate $R^2$ only modestly. DATect's main methodological contribution is therefore not a novel learner, but a leak-aware way to combine per-site modeling, sparse-data feature engineering, and raw-only evaluation.

This perspective also helps interpret the ablations. Interpolated training provides a modest gain, while observation-order lags and derived features show little standalone marginal effect at the one-week horizon---their signal overlaps with persistence and rolling-history terms, and tree models can capture simple nonlinear structure without handcrafted interactions.

\subsection{Spike Detection and Management Relevance}
\label{sec:discussion-spike}

The spike results reveal a different operational problem from regression accuracy. Thresholding the regression ensemble at \SI{20}{\microgrampeg} catches only 12.4\% of held-out below-to-above-threshold transitions even though the model often predicts elevated pre-spike values. This is a predictable consequence of MSE-optimized regression on an imbalanced target: the model is rewarded for shrinking extreme values toward the center of the distribution. For closure-relevant early warning, that conservative bias is a feature for average error but a liability for threshold crossings.

The transition-focused classifier addresses that failure mode by changing the task definition, relying on environmental precursors and gradual toxin accumulation rather than simple persistence (Section~\ref{sec:results-spike-classifier}). The hybrid alert yields 3.1$\times$ higher transition recall than na\"ive persistence, with 85.9\% event recall, but at the cost of substantially lower precision (0.344 vs.\ 0.546) and a lower F2 score (0.661 vs.\ 0.700). This tradeoff is appropriate for a decision-support system: a false alert increases sampling attention, whereas a missed transition can leave managers reacting only after DA has crossed the regulatory limit.

That distinction matters for management value. Quantitative one-week forecasts can support precautionary advisories, sampling prioritization, and more informed timing of openings and closures. The existing PNW HAB Bulletin already provides measurable value through qualitative expert advisories \cite{Jin2024}; DATect complements that workflow with explicit numerical forecasts and a transition-focused alert channel, delivered through a publicly accessible web dashboard (Appendix~\ref{app:webapp}) that supports both real-time forecasting and retrospective analysis across all 10 sites.

\subsection{Limitations}
\label{sec:discussion-limitations}

\begin{enumerate}[nosep]
    \item \textbf{Low-skill Oregon sites}: Three of 10 sites (Cannon Beach, Coos Bay, Newport) have negative $R^2$, and two additional Oregon sites have near-zero skill. Lag-1 autocorrelation analysis (Table~\ref{tab:predictability-ceiling}) reveals persistence-forecast $R^2$ ceilings of 0.003--0.036 for Newport, Coos Bay, and Gold Beach, indicating that consecutive DA measurements at these sites carry very little information about the next observation. This does not preclude improvement from environmental predictors---Gold Beach already exceeds its persistence ceiling ($R^2 = 0.051 > \rho^2 = 0.002$) via SST and BEUTI features---but the low autocorrelation structure means that persistence-based signal contributes little at these sites. The paper cannot fully separate data sparsity from true regional oceanographic complexity; denser Oregon monitoring records and richer environmental predictors may yet reveal exploitable structure currently obscured by irregular sampling.
    \item \textbf{High-risk category weakness}: The High and Extreme DA categories remain difficult, with F1 scores of 0.445 and 0.440, and 13.9\% of Extreme events are misclassified as Low under threshold-based four-category classification.
    \item \textbf{Manual per-site tuning}: Feature subsets, clipping rules, and model choices were selected on a development split. A perturbation study (Appendix~\ref{app:perturbation-sensitivity}) confirms that RF hyperparameters, feature subsets, and monotonic constraints are insensitive to the specific values chosen, while clipping thresholds and per-site model selection are well-calibrated. The held-out seed-123 evaluation and five-seed stability analysis reduce optimistic bias, though a fully nested temporal tuning workflow would provide a still stronger generalization estimate.
    \item \textbf{Uncertainty summaries are not full probabilistic guarantees}: The quantile bands used for forecasts are hybrid empirical intervals, and the bootstrap confidence intervals reported for retrospective metrics resample held-out prediction rows. They are useful uncertainty summaries, but they do not fully resolve temporal dependence or establish exact calibration.
    \item \textbf{Operational deployment remains simplified}: The current system relies on pre-computed cache and a seven-day forecast horizon aligned with the PNW HAB Bulletin. Real-time ingestion and longer forecast horizons remain open engineering and scientific challenges.
    \item \textbf{Sensor continuity}: MODIS-Aqua is approaching end of life, so future operational use will require careful recalibration to successor instruments such as PACE.
\end{enumerate}

\subsection{Future Work}
\label{sec:discussion-future}

The highest-impact next step is improved Oregon monitoring. If data quality is the binding constraint, then denser and more regular sampling should produce larger gains than architecture changes alone. Transfer learning or hierarchical models could help in the interim, but they are unlikely to substitute fully for missing observations.

Additional gains may come from richer observing systems. PACE hyperspectral ocean color \cite{Werdell2019} could improve satellite characterization of bloom conditions, while offshore molecular measurements \cite{Brunson2024} could provide more direct evidence of toxigenic \textit{Pseudo-nitzschia} activity. Sequence models or spatial along-coast models may also become more attractive once the observational record is dense enough to support them.


% ══════════════════════════════════════════════════════════════════════════════
%  6. CONCLUSIONS
% ══════════════════════════════════════════════════════════════════════════════
\section{Conclusions}
\label{sec:conclusions}

DATect provides the first quantitative one-week-ahead domoic acid forecasting system for Washington and Oregon. Under leak-free retrospective evaluation on a held-out test set, all five Washington sites achieve high skill ($R^2 = 0.480$--$0.789$), while Oregon performance is limited by sparser and less regular monitoring. The resulting scientific message is not simply that DATect can forecast DA, but that regional monitoring quality currently places the strongest upper bound on forecast skill.

The paper's main operational result is that magnitude forecasting and closure-warning are related but distinct tasks. Thresholding the regression ensemble at \SI{20}{\microgrampeg} detects only 12.4\% of held-out below-to-above-threshold transitions, whereas the hybrid alert (transition-focused classifier combined with a \SI{12}{\microgrampeg} regression trigger) detects 73.4\% of transitions and 85.9\% of spike events overall. DATect therefore contributes both a quantitative forecasting pipeline and a transition-focused early-warning layer that can complement the qualitative PNW HAB Bulletin and provide a reusable template for other sparse, irregular environmental forecasting problems.


% ══════════════════════════════════════════════════════════════════════════════
%  ACKNOWLEDGMENTS
% ══════════════════════════════════════════════════════════════════════════════
\vspace{6pt}

\authorcontributions{Conceptualization, A.C.; methodology, A.C.; software, A.C.; validation, A.C.; formal analysis, A.C.; investigation, A.C.; data curation, A.C.; writing---original draft preparation, A.C.; writing---review and editing, A.C.\ and J.M.; visualization, A.C.; supervision, J.M. All authors have read and agreed to the published version of the manuscript.}

\funding{This research received no external funding.}

\institutionalreview{Not applicable.}

\informedconsent{Not applicable.}

\dataavailability{The DATect source code and web dashboard are available at \url{https://github.com/ansoncchen/DATect-Forecasting-Domoic-Acid}. Processed datasets and pre-computed cache files are available upon reasonable request.}

\acknowledgments{DA measurement data were provided by the Washington Department of Health, Washington Department of Fish and Wildlife, Oregon Department of Fish and Wildlife, and the Olympic Region Harmful Algal Bloom (ORHAB) Partnership. Satellite data were obtained from the NOAA CoastWatch ERDDAP server. Climate index data were obtained from NOAA ERDDAP. Columbia River discharge data were obtained from the U.S.\ Geological Survey. Computational resources were provided by the University of Washington Hyak high-performance computing cluster.}

\conflictsofinterest{The authors declare no conflict of interest.}


% ══════════════════════════════════════════════════════════════════════════════
%  APPENDICES
% ══════════════════════════════════════════════════════════════════════════════
\appendix

\section{Per-Site Model Configuration}
\label{app:model-config}

Table~\ref{tab:ensemble-config} summarizes the site-wise model selection, clipping rules, and development-set performance tiers used during configuration selection.

\begin{table}[H]
\caption{Per-site ensemble configuration showing model selection, prediction clipping, and development-set performance tier used during configuration selection (see Table~\ref{tab:per-site-results} for held-out test-set $R^2$).}
\label{tab:ensemble-config}
\small
\begin{tabular}{lccccl}
\toprule
\textbf{Site} & $\boldsymbol{w_{\text{XGB}}}$ & $\boldsymbol{w_{\text{RF}}}$ & \textbf{Clip $q$} & \textbf{Clip max} & \textbf{Tier} \\
\midrule
Copalis       & 0.00 & 1.00 & 0.97 & --- & High-skill \\
Twin Harbors  & 0.00 & 1.00 & 0.98 & --- & High-skill \\
Quinault      & 0.00 & 1.00 & 0.98 & --- & High-skill \\
Long Beach    & 1.00 & 0.00 & 0.98 & --- & High-skill \\
Kalaloch      & 0.00 & 1.00 & 0.95 & 80 & High-skill \\
\addlinespace
Clatsop Beach & 1.00 & 0.00 & --- & --- & Moderate \\
Gold Beach    & 1.00 & 0.00 & 0.95 & --- & Moderate \\
\addlinespace
Cannon Beach  & 0.00 & 1.00 & 0.95 & 80 & Low-skill \\
\addlinespace
Coos Bay      & 0.00 & 1.00 & 0.97 & --- & Low-skill \\
Newport       & 1.00 & 0.00 & 0.98 & --- & Low-skill \\
\bottomrule
\end{tabular}
\vspace{0.3em}\\
\footnotesize{Clip~$q$: predictions clipped to this quantile of training targets. ``---'' = global default ($q = 0.99$) or no hard ceiling. Clip max in \si{\microgrampeg}. All sites converged to single-model selection ($w \in \{0,1\}$); blended weights were tested but did not improve over single-model selection (Appendix~\ref{app:weight-robustness}).}
\end{table}

\section{Classification Details}
\label{app:classification-details}

\begin{table}[H]
\caption{DA risk categories used in DATect classification.}
\label{tab:da-categories}
\small
\begin{tabular}{lcl}
\toprule
\textbf{Category} & \textbf{DA Range (\si{\microgrampeg})} & \textbf{Interpretation} \\
\midrule
Low      & 0--5   & Safe for harvest \\
Moderate & 5--20  & Monitor closely \\
High     & 20--40 & Elevated toxicity; closures likely \\
Extreme  & $>$40  & Dangerous; immediate closure \\
\bottomrule
\end{tabular}
\end{table}

\begin{table}[H]
\caption{Per-category classification performance on the held-out test set ($n = 2{,}181$). The Low category achieves the highest reliability, while High and Extreme categories are limited by sample size. Overall accuracy 75.1\% (macro-average precision 0.605, recall 0.581, F1 0.591).}
\label{tab:classification}
\small
\begin{tabular}{lrrrrr}
\toprule
\textbf{Category} & \textbf{DA Range} & $\boldsymbol{N}$ & \textbf{Precision} & \textbf{Recall} & \textbf{F1} \\
\midrule
Low      & 0--5 \si{\microgrampeg}  & 1{,}382 & 0.865 & 0.858 & 0.862 \\
Moderate & 5--20 \si{\microgrampeg} & 523 & 0.593 & 0.642 & 0.617 \\
High     & 20--40 \si{\microgrampeg} & 168 & 0.456 & 0.435 & 0.445 \\
Extreme  & $\geq$40 \si{\microgrampeg}  & 108 & 0.506 & 0.389 & 0.440 \\
\midrule
\multicolumn{2}{l}{\textbf{Overall (macro)}} & \textbf{2,181} & \textbf{0.605} & \textbf{0.581} & \textbf{0.591} \\
\multicolumn{2}{l}{Overall accuracy} & & \multicolumn{3}{r}{75.1\%} \\
\bottomrule
\end{tabular}
\end{table}

\begin{table}[H]
\caption{Confusion matrix for four-category DA classification. Most errors occur between adjacent categories.}
\label{tab:confusion}
\small
\begin{tabular}{l|rrrr}
\toprule
\textbf{Actual \textbackslash\ Predicted} & \textbf{Low} & \textbf{Moderate} & \textbf{High} & \textbf{Extreme} \\
\midrule
Low      & 1183 & 142 & 27 & 30 \\
Moderate & 164  & 327 & 30 & 2 \\
High     & 16   & 80  & 67 & 5 \\
Extreme  & 15   & 22  & 32 & 39 \\
\bottomrule
\end{tabular}
\end{table}

\section{Component Ablation}
\label{app:component-ablation}

\begin{table}[H]
\caption{Component ablation study showing the contribution of each DATect design choice. Each row removes one component from the full system. $\Delta R^2$ and $\Delta$MAE are relative to the baseline. Evaluated on the development set (seed = 42, $n = 1{,}177$); the held-out test set (seed = 123) was reserved for final reporting.}
\label{tab:ablation}
\small
\begin{tabular}{lcccl}
\toprule
\textbf{Configuration} & $\boldsymbol{R^2}$ & $\boldsymbol{\Delta R^2}$ & \textbf{MAE} & \textbf{Notes} \\
\midrule
Full DATect (baseline)        & 0.414 & ---      & 6.38 & Current configuration (dev set) \\
No per-site customization     & 0.256 & $-0.158$ & 6.69 & Largest impact \\
No interpolated training      & 0.385 & $-0.029$ & 6.61 & Data augmentation effect \\
No observation-order lags     & 0.413 & $-0.001$ & 6.39 & Minimal standalone impact \\
No derived features           & 0.415 & $+0.001$ & 6.38 & Negligible impact \\
\bottomrule
\end{tabular}
\end{table}

\section{Per-Site Hyperparameter Configuration}
\label{app:hyperparams}

Table~\ref{tab:app-xgb} presents the full XGBoost hyperparameter configuration for each site.

\begin{table}[H]
\caption{Per-site XGBoost hyperparameter overrides. Parameters not listed use global defaults ($n_{\text{est}} = 400$, depth $= 6$, lr $= 0.05$, subsample $= 0.85$, colsample\_bytree $= 0.85$, colsample\_bylevel $= 0.8$, $\alpha = 0.1$, $\lambda = 1.0$, $\gamma = 0.1$, min\_child\_wt $= 3$, tree\_method $=$ hist).}
\label{tab:app-xgb}
\small
\begin{tabular}{lccccccccc}
\toprule
\textbf{Site} & \textbf{depth} & $\boldsymbol{n_{\text{est}}}$ & \textbf{lr} & \textbf{mcw} & $\boldsymbol{\alpha}$ & $\boldsymbol{\lambda}$ & $\boldsymbol{\gamma}$ & \textbf{ss} & \textbf{cs} \\
\midrule
Kalaloch      & 2 & 80  & 0.02 & 12 & 2.0 & 10.0 & 2.0 & 0.6 & 0.6 \\
Quinault      & 3 & 200 & 0.03 & 7  & 0.3 & 2.0 & 0.3 & 0.8 & 0.8 \\
Copalis       & 2 & 100 & 0.03 & 10 & 1.0 & 5.0 & 1.0 & 0.7 & 0.7 \\
Twin Harbors  & 3 & 150 & 0.03 & 8  & 0.5 & 3.0 & 0.5 & 0.8 & 0.8 \\
Long Beach    & 3 & 250 & 0.03 & 7  & 0.3 & 2.0 & 0.3 & 0.8 & 0.8 \\
Clatsop Beach & \multicolumn{9}{c}{\textit{(global defaults)}} \\
Cannon Beach  & 2 & 80  & 0.02 & 12 & 1.5 & 7.0  & 1.5 & 0.6 & 0.6 \\
Newport       & 2 & 100 & 0.02 & 15 & 2.0 & 10.0 & 2.0 & 0.6 & 0.6 \\
Coos Bay      & 2 & 80  & 0.02 & 15 & 2.0 & 10.0 & 2.0 & 0.6 & 0.6 \\
Gold Beach    & 2 & 100 & 0.02 & 15 & 2.0 & 10.0 & 2.0 & 0.6 & 0.6 \\
\bottomrule
\end{tabular}
\vspace{0.3em}\\
\footnotesize{mcw = min\_child\_weight; ss = subsample; cs = colsample\_bytree. Clatsop Beach uses all global defaults. All other sites override every listed parameter.}
\end{table}

Sites with low historical skill (Cannon Beach, Coos Bay, Gold Beach, Newport) use the conservative Random Forest configuration: $n_{\text{est}} = 200$, max depth $= 6$, min\_samples\_split $= 10$, min\_samples\_leaf $= 5$, max\_features $= 0.5$.

\section{Feature Subsets by Site}
\label{app:features}

Table~\ref{tab:app-features} summarizes the feature categories included at each site. A checkmark indicates the category is included in that site's curated feature subset.

\begin{table}[H]
\caption{Feature category inclusion by site. Sites with ``All'' use the full feature set without subsetting. ``partial'' indicates a subset of the group (e.g., only \texttt{modis-sst} and \texttt{pdo} from Env.\ Core). Newport, Coos Bay, and Gold Beach use a single lag feature (\texttt{da\_raw\_prev\_obs\_1}) rather than a complete lag group; Cannon Beach uses the first two obs lags only (\texttt{obs\_1}, \texttt{obs\_2}). These sites were trimmed to $\leq$7 features to reduce overfitting given small training $N$. Features explored but removed from the pipeline are listed in Appendix~\ref{app:feature-engineering}.}
\label{tab:app-features}
\footnotesize
\begin{tabular}{l|ccccccccc}
\toprule
\textbf{Site} & \rotatebox{70}{Persistence} & \rotatebox{70}{Full Lags} & \rotatebox{70}{Short Lags} & \rotatebox{70}{Full Rolling} & \rotatebox{70}{Short Rolling} & \rotatebox{70}{Env.\ Core} & \rotatebox{70}{Climate} & \rotatebox{70}{Discharge} & \rotatebox{70}{PN} \\
\midrule
Kalaloch      & $\checkmark$ & & $\checkmark$ & & & & $\checkmark$ & $\checkmark$ & $\checkmark$ \\
Quinault      & $\checkmark$ & $\checkmark$ & & & $\checkmark$ & $\checkmark$ & $\checkmark$ & $\checkmark$ & \\
Copalis       & $\checkmark$ & & $\checkmark$ & & $\checkmark$ & & $\checkmark$ & $\checkmark$ & \\
Twin Harbors  & $\checkmark$ & $\checkmark$ & & & $\checkmark$ & partial & $\checkmark$ & $\checkmark$ & \\
Long Beach    & $\checkmark$ & $\checkmark$ & & $\checkmark$ & & $\checkmark$ & $\checkmark$ & $\checkmark$ & \\
Clatsop Beach & \multicolumn{9}{c}{\textit{All (no subset)}} \\
Cannon Beach  & $\checkmark$ & & $\checkmark$ & & & partial & & & \\
Newport       & $\checkmark$ & & & & & partial & & & \\
Coos Bay      & $\checkmark$ & & & & & partial & & & \\
Gold Beach    & $\checkmark$ & & & & & partial & & & \\
\bottomrule
\end{tabular}
\end{table}

\section{Features Explored but Removed}
\label{app:dropped-features}
\label{app:derived-features}
\label{app:feature-engineering}

During development, we explored a number of additional features motivated by domain knowledge of harmful algal bloom dynamics. Individual feature ablation (Table~\ref{tab:feature-ablation}) confirmed that none contributed more than $|\Delta R^2| = 0.005$ to ensemble performance, and all have been removed from the final pipeline. We document them here for the benefit of future researchers.

\subsection*{Derived Environmental Features}

\begin{itemize}[nosep]
    \item \textbf{Marine heatwave flag}: Binary indicator for SST anomaly $> \SI{1.5}{\degreeCelsius}$, motivated by the 2015 ``Blob'' event \cite{Bond2015, McCabe2016}. $\Delta R^2 = -0.004$.
    \item \textbf{BEUTI nonlinear terms}: $\text{BEUTI}^2$ for the upwelling ``Goldilocks zone'' \cite{Trainer2020} and a relaxation indicator for onshore transport \cite{Trainer2002}. $\Delta R^2 = -0.001$ and $-0.002$.
    \item \textbf{PDO--ONI phase coherence}: Binary flag for co-occurrence of positive PDO and El Ni\~no \cite{McKibben2017}. $\Delta R^2 = -0.002$.
    \item \textbf{\textit{Pseudo-nitzschia} threshold}: Binary indicator for cell counts $> 50{,}000$~\si{\cellsperL}. $\Delta R^2 = -0.005$.
    \item \textbf{Fluorescence efficiency}: FLH\,/\,Chl-$a$ ratio as a physiological stress proxy. $\Delta R^2 = -0.001$.
    \item \textbf{K490 quadratic}: $\text{K490}^2$ for turbidity nonlinearity. $\Delta R^2 = +0.001$ (removing \textit{improves} performance).
\end{itemize}

\subsection*{Satellite Observations}

Several satellite products were downloaded and tested but contributed negligibly:
\begin{itemize}[nosep]
    \item \textbf{MODIS chlorophyll-$a$} and \textbf{chlorophyll-$a$ anomaly}: $\Delta R^2 = -0.004$ and $-0.001$, respectively.
    \item \textbf{MODIS K490} (diffuse attenuation coefficient): $\Delta R^2 = +0.001$ (removing improves performance).
    \item \textbf{MODIS PAR} (photosynthetically active radiation): $< 1\%$ feature importance.
\end{itemize}
Removing all six satellite-derived features simultaneously costs only $\Delta R^2 = -0.003$.

\subsection*{Temporal and Other Features}

Additional temporal encodings (\texttt{quarter}, \texttt{is\_bloom\_season}, sine/cosine of week-of-year, \texttt{days\_since\_start}), a recency feature (\texttt{weeks\_since\_last\_raw}), site coordinates (\texttt{lat}, \texttt{lon}), and rolling-mean features for 8- and 12-week windows all had $< 1\%$ feature importance and were removed.


\section{Web Application Details}
\label{app:webapp}

\subsection*{Architecture and Deployment}

\begin{figure}[H]
\centering
\includegraphics[width=\textwidth]{figures/fig3_architecture.png}
\caption{DATect system architecture. Data from satellite, climate, hydrological, and biological sources are processed into a unified dataset. Raw DA measurements are also read directly by the forecasting engine for raw-only persistence features. Results are served through a FastAPI backend to a React dashboard, with a pre-computed cache accelerating retrospective evaluations.}
\label{fig:architecture}
\end{figure}

The DATect backend serves forecast requests, retrospective evaluation results, and visualization data (correlation heatmaps, sensitivity analyses, spectral analyses, feature importance rankings). A pre-computed cache system stores results from computationally intensive retrospective evaluations as JSON and Parquet files, with an optional Redis backend providing 100$\times$ faster cache reads for production deployments. The retrospective evaluation cache---encompassing all 2,181 validation points across 10 sites and 8 model configurations---was pre-computed on the University of Washington Hyak high-performance computing cluster (\textit{klone} partition, STF allocation). Because each test point requires training a fresh ensemble on its own chronological training subset, the per-anchor computation is embarrassingly parallel and well-suited to HPC execution; running the full retrospective evaluation sequentially on a laptop would require several days, whereas the Hyak cluster completes the computation in approximately 2--3 hours across parallel jobs.

The system is containerized with Docker and deployed on Google Cloud Run, a serverless platform that scales to zero when idle and scales up to three concurrent instances under load. Cold start time is approximately 10--15 seconds. The deployment pipeline uses Google Cloud Build for continuous integration.

\subsection*{Operational Modes}

The dashboard supports two modes. In \textbf{real-time forecasting mode}, users select a target date and monitoring site; DATect trains a fresh model on all data up to the anchor date, returning a point prediction, 90\% confidence interval, risk category, and per-category probability distribution. In \textbf{retrospective analysis mode}, users select a model type and task to browse pre-computed validation results across all 2,181 test points, with interactive Plotly charts for scatter plots, confusion matrices, residual distributions, and time series overlays, filterable by site, date range, and DA risk category. Seven additional visualization types---including correlation heatmaps, sensitivity analyses, spectral density plots, SHAP-style feature waterfalls, and an interactive geographic map---are available for exploratory analysis. We omit a dashboard screenshot here because the manuscript no longer includes placeholder figures; the live interface is available through the deployed application and local development workflow described in the repository.


\section{Individual Feature Ablation}
\label{app:feature-ablation}

Table~\ref{tab:feature-ablation} presents a fine-grained ablation in which each row removes one feature or feature group from the full system. All changes fall within sampling variability ($|\Delta R^2| \leq 0.005$), confirming negligible marginal contributions. All features listed were removed from the final pipeline; see Appendix~\ref{app:dropped-features} for documentation.

\begin{table}[H]
\caption{Feature group ablation. Each row removes one feature or group from the full system. All changes fall within sampling variability ($|\Delta R^2| \leq 0.005$), confirming negligible marginal contributions. Evaluated on the development set (seed = 42) under an earlier per-site configuration (baseline $R^2 = 0.457$); the component ablation table (Table~\ref{tab:ablation}) uses the final configuration ($R^2 = 0.414$). Since all $|\Delta R^2| \leq 0.005$, this conclusion is robust to the configuration difference.}
\label{tab:feature-ablation}
\small
\begin{tabular}{lcccc}
\toprule
\textbf{Feature(s) removed} & $\boldsymbol{R^2}$ & $\boldsymbol{\Delta R^2}$ & \textbf{MAE} & $\boldsymbol{\Delta}$\textbf{MAE} \\
\midrule
Full DATect (baseline) & 0.457 & --- & 6.15 & --- \\
\addlinespace
All satellite-derived (6 features) & 0.454 & $-0.003$ & 6.19 & $+0.03$ \\
All derived env.\ features (4 features) & 0.455 & $-0.002$ & 6.18 & $+0.03$ \\
PN features (2) & 0.452 & $-0.005$ & 6.22 & $+0.06$ \\
Rolling std features (3) & 0.457 & $-0.001$ & 6.16 & $+0.00$ \\
\addlinespace
All low-importance combined (8 features)$^\dagger$ & 0.453 & $-0.004$ & 6.19 & $+0.04$ \\
\bottomrule
\end{tabular}
\vspace{0.3em}\\
\footnotesize{$^\dagger$Includes: K490, K490$^2$, fluorescence efficiency, Chl-$a$ anomaly, PDO--ONI phase, BEUTI relaxation, PN threshold, and Chl-$a$. Individual feature-level removals (not shown) all fall within $|\Delta R^2| \leq 0.005$.}
\end{table}


\section{Feature Importance Analysis}
\label{app:feature-importance}

Figure~\ref{fig:feature-importance} presents mean feature importance (gain) across all 1,177 development-set predictions. Persistence features (\texttt{last\_observed\_da\_raw}, \texttt{weeks\_since\_last\_spike}) dominate, accounting for over half of total model gain and confirming the strong autocorrelation in DA time series. Rolling statistics and observation-order lag features contribute meaningful secondary signal, while environmental features (PDO, SST, ONI) appear in the top 15 but at substantially lower magnitude. Importantly, high feature importance does not imply large marginal contribution: lag features rank highly because the model uses them frequently, but the component ablation (Table~\ref{tab:ablation}) shows that removing them costs only $\Delta R^2 = -0.001$, indicating overlap with persistence and rolling-history terms.

\begin{figure}[H]
\centering
\includegraphics[width=0.85\textwidth]{figures/fig6_feature_importance.png}
\caption{Mean feature importance (gain) across all development-set predictions ($n = 1{,}177$), averaged over per-site ensemble models. Features are color-coded by category. Persistence features dominate, consistent with strong DA autocorrelation at weekly timescales; environmental features contribute meaningful but secondary signal.}
\label{fig:feature-importance}
\end{figure}


\section{Model Selection Robustness}
\label{app:weight-robustness}

Since DATect assigns each site a single model (XGBoost or Random Forest) based on development-set performance, weight robustness reduces to model-selection robustness: how sensitive are aggregate results to the per-site XGB vs.\ RF choice? Table~\ref{tab:weight-robustness} compares three strategies evaluated on the held-out test set ($n = 2{,}181$, seed = 123).

\begin{table}[H]
\caption{Model selection robustness on the held-out test set ($n = 2{,}181$, seed = 123). The oracle selects the better single model per site using test-set $R^2$ directly. Bootstrap 95\% CIs computed with $B = 10{,}000$ resamples.}
\label{tab:weight-robustness}
\small
\begin{tabular}{lccc}
\toprule
\textbf{Strategy} & \textbf{$R^2$} & \textbf{MAE (\si{\microgrampeg})} & \textbf{RMSE (\si{\microgrampeg})} \\
\midrule
XGBoost at all sites                    & 0.221 & 6.38 & 17.72 \\
RF at all sites                         & 0.190 & 6.47 & 18.07 \\
Per-site dev-set selection (this paper) & 0.215 & 6.42 & 17.79 \\
\bottomrule
\end{tabular}
{\footnotesize The oracle---selecting the better single model per site using test-set $R^2$---achieves pooled $R^2 = 0.210$, MAE $= 6.41$~\si{\microgrampeg}. This is slightly \textit{lower} than the dev-set selection ($\Delta R^2 = -0.005$) because pooled $R^2$ is not a weighted average of per-site values: the oracle's per-site-optimal choices can concentrate residual error at high-$N$ sites that dominate the pooled metric. The near-equivalence of all four strategies ($R^2 = 0.190$--$0.221$) confirms that model-family selection is not a meaningful source of variance.}
\end{table}


\section{Perturbation Sensitivity Analysis}
\label{app:perturbation-sensitivity}

To validate that per-site design choices are not fragile artifacts of the seed-42 development set, we conduct a systematic perturbation study. Each experiment modifies a single design dimension while holding all others fixed, evaluated on the development set ($n = 1{,}177$, seed = 42, 20\% sample). Table~\ref{tab:perturbation} reports the results.

\begin{table}[H]
\caption{Perturbation sensitivity analysis. Each row modifies one design choice relative to the baseline configuration. $\Delta R^2$ is relative to baseline ($R^2 = 0.409$). The baseline here differs slightly from the component ablation table ($R^2 = 0.414$, Table~\ref{tab:ablation}) because the perturbation study was run on an intermediate configuration checkpoint; both use the same development set (seed = 42, $n = 1{,}177$).}
\label{tab:perturbation}
\small
\begin{tabular}{lrrr}
\toprule
\textbf{Perturbation} & $\boldsymbol{R^2}$ & $\boldsymbol{\Delta R^2}$ & \textbf{MAE} \\
\midrule
\multicolumn{4}{l}{\textit{Model selection}} \\
\quad Swap all RF-only sites $\rightarrow$ XGB & 0.381 & $-0.029$ & 6.47 \\
\quad Swap all XGB-only sites $\rightarrow$ RF & 0.405 & $-0.005$ & 6.38 \\
\quad No per-site configuration               & 0.307 & $-0.103$ & 6.66 \\
\midrule
\multicolumn{4}{l}{\textit{Random Forest hyperparameters}} \\
\quad Shallow (depth = 6, $n_{\text{est}}$ = 200)   & 0.410 & $+0.000$ & 6.40 \\
\quad Deep (depth = 16, $n_{\text{est}}$ = 600)      & 0.409 & $-0.000$ & 6.42 \\
\quad More trees ($n_{\text{est}}$ = 800)             & 0.410 & $+0.000$ & 6.42 \\
\quad Less regularized (leaf = 1)                     & 0.409 & $-0.001$ & 6.43 \\
\midrule
\multicolumn{4}{l}{\textit{Feature and clipping choices}} \\
\quad All features at all sites                       & 0.415 & $+0.006$ & 6.45 \\
\quad Minimal features (4 only)                       & 0.412 & $+0.002$ & 6.66 \\
\quad Relax clipping (quantile = 0.99)                & 0.380 & $-0.029$ & 6.49 \\
\quad Disable clipping                                & 0.384 & $-0.026$ & 6.39 \\
\quad Disable monotonic constraints                   & 0.410 & $+0.000$ & 6.43 \\
\bottomrule
\end{tabular}
\end{table}

Three findings emerge. First, Random Forest hyperparameters are genuinely insensitive: across four configurations spanning depth 6--16 and 200--800 estimators, no perturbation exceeds $|\Delta R^2| = 0.001$. This validates the claim in Section~\ref{sec:results-ablation} that RF robustness to hyperparameters does not require per-site RF tuning. Second, per-site prediction clipping provides meaningful regularization: relaxing all per-site quantile thresholds from their current values (0.95--0.98) to a uniform 0.99 costs $\Delta R^2 = -0.029$; notably, disabling clipping entirely costs only $-0.026$, revealing that the harm is not from clipping itself but from replacing well-calibrated site-specific thresholds with an overly permissive uniform one. Third, feature subsets have negligible pooled impact ($+0.006$ with all features, $+0.002$ with only four persistence features), confirming that predictive skill is dominated by the DA persistence signal rather than feature engineering.

Multi-seed evaluation further characterizes the stability of these results. Table~\ref{tab:seed-stability} reports per-site $R^2$ across five random seeds with per-site configuration enabled.

\begin{table}[H]
\caption{Per-site $R^2$ across five random seeds (20\% sample, per-site configuration enabled). Sites are grouped by region.}
\label{tab:seed-stability}
\small
\begin{tabular}{lrrrrrrr}
\toprule
\textbf{Site} & \textbf{Seed 42} & \textbf{123} & \textbf{456} & \textbf{789} & \textbf{1337} & \textbf{Mean} & \textbf{Std} \\
\midrule
\multicolumn{8}{l}{\textit{Washington}} \\
Copalis       & 0.752 & 0.764 & 0.769 & 0.712 & 0.733 & 0.746 & 0.021 \\
Quinault      & 0.491 & 0.584 & 0.677 & 0.680 & 0.519 & 0.590 & 0.078 \\
Twin Harbors  & 0.514 & 0.566 & 0.579 & 0.709 & 0.566 & 0.587 & 0.065 \\
Long Beach    & 0.643 & 0.526 & 0.605 & 0.642 & 0.516 & 0.586 & 0.055 \\
Kalaloch      & 0.618 & 0.627 & 0.454 & 0.628 & 0.537 & 0.573 & 0.069 \\
\midrule
\multicolumn{8}{l}{\textit{Oregon}} \\
Clatsop Beach & 0.374 & 0.305 & $-0.238$ & 0.666 & 0.296 & 0.280 & 0.292 \\
Coos Bay      & 0.358 & $-0.004$ & 0.095 & $-0.001$ & 0.153 & 0.120 & 0.133 \\
Cannon Beach  & 0.014 & $-0.071$ & 0.138 & $-0.055$ & $-0.071$ & $-0.009$ & 0.080 \\
Gold Beach    & $-0.053$ & 0.035 & $-0.251$ & 0.227 & 0.006 & $-0.007$ & 0.154 \\
Newport       & 0.119 & $-0.163$ & $-0.950$ & $-0.113$ & $-0.226$ & $-0.267$ & 0.361 \\
\midrule
\multicolumn{2}{l}{Pooled} & & & & & 0.265 & 0.102 \\
\bottomrule
\end{tabular}
\end{table}

Washington sites show stable performance across seeds (mean $R^2 = 0.57$--$0.75$, std $\leq 0.08$), while Oregon sites exhibit high variance (std up to 0.36 for Newport). The regional skill divide persists across all five seeds, confirming that it reflects fundamental data quality differences rather than seed-specific tuning artifacts.


% ══════════════════════════════════════════════════════════════════════════════
%  REFERENCES
% ══════════════════════════════════════════════════════════════════════════════
\externalbibliography{yes}
\bibliography{references}

\end{document}
